\documentclass[a4paper,12pt,twoside]{memoir}

% Castellano
\usepackage[spanish,es-tabla]{babel}
\selectlanguage{spanish}
\usepackage[utf8]{inputenc}
\usepackage[T1]{fontenc}
\usepackage{lmodern} % scalable font
\usepackage{microtype}
\usepackage{placeins}

\RequirePackage{booktabs}
\RequirePackage[table]{xcolor}
\RequirePackage{xtab}
\RequirePackage{multirow}

% Links
\PassOptionsToPackage{hyphens}{url}\usepackage[colorlinks]{hyperref}
\hypersetup{
	allcolors = {red}
}

% Ecuaciones
\usepackage{amsmath}

% Rutas de fichero / paquete
\newcommand{\ruta}[1]{{\sffamily #1}}

% Párrafos
\nonzeroparskip

% Huérfanas y viudas
\widowpenalty100000
\clubpenalty100000

% Evitar solapes en el header
\nouppercaseheads

% Imagenes
\usepackage{graphicx}
\newcommand{\imagen}[2]{
	\begin{figure}[!h]
		\centering
		\includegraphics[width=0.9\textwidth]{#1}
		\caption{#2}\label{fig:#1}
	\end{figure}
	\FloatBarrier
}

\newcommand{\imagenflotante}[2]{
	\begin{figure}%[!h]
		\centering
		\includegraphics[width=0.9\textwidth]{#1}
		\caption{#2}\label{fig:#1}
	\end{figure}
}



% El comando \figura nos permite insertar figuras comodamente, y utilizando
% siempre el mismo formato. Los parametros son:
% 1 -> Porcentaje del ancho de página que ocupará la figura (de 0 a 1)
% 2 --> Fichero de la imagen
% 3 --> Texto a pie de imagen
% 4 --> Etiqueta (label) para referencias
% 5 --> Opciones que queramos pasarle al \includegraphics
% 6 --> Opciones de posicionamiento a pasarle a \begin{figure}
\newcommand{\figuraConPosicion}[6]{%
  \setlength{\anchoFloat}{#1\textwidth}%
  \addtolength{\anchoFloat}{-4\fboxsep}%
  \setlength{\anchoFigura}{\anchoFloat}%
  \begin{figure}[#6]
    \begin{center}%
      \Ovalbox{%
        \begin{minipage}{\anchoFloat}%
          \begin{center}%
            \includegraphics[width=\anchoFigura,#5]{#2}%
            \caption{#3}%
            \label{#4}%
          \end{center}%
        \end{minipage}
      }%
    \end{center}%
  \end{figure}%
}

%
% Comando para incluir imágenes en formato apaisado (sin marco).
\newcommand{\figuraApaisadaSinMarco}[5]{%
  \begin{figure}%
    \begin{center}%
    \includegraphics[angle=90,height=#1\textheight,#5]{#2}%
    \caption{#3}%
    \label{#4}%
    \end{center}%
  \end{figure}%
}
% Para las tablas
\newcommand{\otoprule}{\midrule [\heavyrulewidth]}
%
% Nuevo comando para tablas pequeñas (menos de una página).
\newcommand{\tablaSmall}[5]{%
 \begin{table}
  \begin{center}
   \rowcolors {2}{gray!35}{}
   \begin{tabular}{#2}
    \toprule
    #4
    \otoprule
    #5
    \bottomrule
   \end{tabular}
   \caption{#1}
   \label{tabla:#3}
  \end{center}
 \end{table}
}

%
%Para el float H de tablaSmallSinColores
\usepackage{float}

%
% Nuevo comando para tablas pequeñas (menos de una página).
\newcommand{\tablaSmallSinColores}[5]{%
 \begin{table}[H]
  \begin{center}
   \begin{tabular}{#2}
    \toprule
    #4
    \otoprule
    #5
    \bottomrule
   \end{tabular}
   \caption{#1}
   \label{tabla:#3}
  \end{center}
 \end{table}
}

\newcommand{\tablaApaisadaSmall}[5]{%
\begin{landscape}
  \begin{table}
   \begin{center}
    \rowcolors {2}{gray!35}{}
    \begin{tabular}{#2}
     \toprule
     #4
     \otoprule
     #5
     \bottomrule
    \end{tabular}
    \caption{#1}
    \label{tabla:#3}
   \end{center}
  \end{table}
\end{landscape}
}

%
% Nuevo comando para tablas grandes con cabecera y filas alternas coloreadas en gris.
\newcommand{\tabla}[6]{%
  \begin{center}
    \tablefirsthead{
      \toprule
      #5
      \otoprule
    }
    \tablehead{
      \multicolumn{#3}{l}{\small\sl continúa desde la página anterior}\\
      \toprule
      #5
      \otoprule
    }
    \tabletail{
      \hline
      \multicolumn{#3}{r}{\small\sl continúa en la página siguiente}\\
    }
    \tablelasttail{
      \hline
    }
    \bottomcaption{#1}
    \rowcolors {2}{gray!35}{}
    \begin{xtabular}{#2}
      #6
      \bottomrule
    \end{xtabular}
    \label{tabla:#4}
  \end{center}
}

%
% Nuevo comando para tablas grandes con cabecera.
\newcommand{\tablaSinColores}[6]{%
  \begin{center}
    \tablefirsthead{
      \toprule
      #5
      \otoprule
    }
    \tablehead{
      \multicolumn{#3}{l}{\small\sl continúa desde la página anterior}\\
      \toprule
      #5
      \otoprule
    }
    \tabletail{
      \hline
      \multicolumn{#3}{r}{\small\sl continúa en la página siguiente}\\
    }
    \tablelasttail{
      \hline
    }
    \bottomcaption{#1}
    \begin{xtabular}{#2}
      #6
      \bottomrule
    \end{xtabular}
    \label{tabla:#4}
  \end{center}
}

%
% Nuevo comando para tablas grandes sin cabecera.
\newcommand{\tablaSinCabecera}[5]{%
  \begin{center}
    \tablefirsthead{
      \toprule
    }
    \tablehead{
      \multicolumn{#3}{l}{\small\sl continúa desde la página anterior}\\
      \hline
    }
    \tabletail{
      \hline
      \multicolumn{#3}{r}{\small\sl continúa en la página siguiente}\\
    }
    \tablelasttail{
      \hline
    }
    \bottomcaption{#1}
  \begin{xtabular}{#2}
    #5
   \bottomrule
  \end{xtabular}
  \label{tabla:#4}
  \end{center}
}



\definecolor{cgoLight}{HTML}{EEEEEE}
\definecolor{cgoExtralight}{HTML}{FFFFFF}

%
% Nuevo comando para tablas grandes sin cabecera.
\newcommand{\tablaSinCabeceraConBandas}[5]{%
  \begin{center}
    \tablefirsthead{
      \toprule
    }
    \tablehead{
      \multicolumn{#3}{l}{\small\sl continúa desde la página anterior}\\
      \hline
    }
    \tabletail{
      \hline
      \multicolumn{#3}{r}{\small\sl continúa en la página siguiente}\\
    }
    \tablelasttail{
      \hline
    }
    \bottomcaption{#1}
    \rowcolors[]{1}{cgoExtralight}{cgoLight}

  \begin{xtabular}{#2}
    #5
   \bottomrule
  \end{xtabular}
  \label{tabla:#4}
  \end{center}
}




\graphicspath{ {./img/} }

% Capítulos
\chapterstyle{bianchi}
\newcommand{\capitulo}[2]{
	\setcounter{chapter}{#1}
	\setcounter{section}{0}
	\setcounter{figure}{0}
	\setcounter{table}{0}
	\chapter*{#2}
	\addcontentsline{toc}{chapter}{#2}
	\markboth{#2}{#2}
}

% Apéndices
\renewcommand{\appendixname}{Apéndice}
\renewcommand*\cftappendixname{\appendixname}

\newcommand{\apendice}[1]{
	%\renewcommand{\thechapter}{A}
	\chapter{#1}
}

\renewcommand*\cftappendixname{\appendixname\ }

% Formato de portada
\makeatletter
\usepackage{xcolor}
\newcommand{\tutor}[1]{\def\@tutor{#1}}
\newcommand{\course}[1]{\def\@course{#1}}
\definecolor{cpardoBox}{HTML}{E6E6FF}
\def\maketitle{
  \null
  \thispagestyle{empty}
  % Cabecera ----------------
\noindent\includegraphics[width=\textwidth]{cabecera}\vspace{1cm}%
  \vfill
  % Título proyecto y escudo informática ----------------
  \colorbox{cpardoBox}{%
    \begin{minipage}{.8\textwidth}
      \vspace{.5cm}\Large
      \begin{center}
      \textbf{TFG del Grado en Ingeniería Informática}\vspace{.6cm}\\
      \textbf{\LARGE\@title{}}
      \end{center}
      \vspace{.2cm}
    \end{minipage}

  }%
  \hfill\begin{minipage}{.20\textwidth}
    \includegraphics[width=\textwidth]{escudoInfor}
  \end{minipage}
  \vfill
  % Datos de alumno, curso y tutores ------------------
  \begin{center}%
  {%
    \noindent\LARGE
    Presentado por \@author{}\\ 
    en Universidad de Burgos --- \@date{}\\
    Tutor: \@tutor{}\\
  }%
  \end{center}%
  \null
  \cleardoublepage
  }
\makeatother


% Datos de portada
\title{título del TFG \\Documentación Técnica}
\author{nombre alumno}
\tutor{nombre tutor}
\date{\today}

\begin{document}

\maketitle



\cleardoublepage



%%%%%%%%%%%%%%%%%%%%%%%%%%%%%%%%%%%%%%%%%%%%%%%%%%%%%%%%%%%%%%%%%%%%%%%%%%%%%%%%%%%%%%%%



\frontmatter


\clearpage

% Indices
\tableofcontents

\clearpage

\listoffigures

\clearpage

\listoftables

\clearpage

\mainmatter

\appendix

\apendice{Plan de Proyecto Software}

\section{Introducción}

En un proyecto con un equipo completo de personas trabajando, podríamos definir la planificación de un proyecto software como una etapa que implica detallar los objetivos y el alcance del proyecto, identificar los requisitos y recursos necesarios, establecer un cronograma con hitos clave, asignar tareas a los miembros del equipo y prever los posibles riesgos. 

Dentro del contexto de este trabajo académico, se va a detallar en qué ha consistido la planificación temporal y se tratará de hacer un estudio de viabilidad realista.


\section{Planificación temporal}

En la memoria se ha indicado que se ha utilizado una metodología ágil de gestión de proyectos, con clara fundamentación en Scrum\citep{wiki:scrum}. Algunos de los aspectos más relevantes que han cubierto esta filosofía han sido:

\begin{itemize}
\tightlist
\item  
Desarrollo incremental a través de iteraciones llamadas \emph{sprints}.
\item
Utilización de repositorio \emph{Git} para solicitud de mejoras y para realizar un seguimiento de la evolución, con acceso para los tutores desde las fases iniciales. 
\item
Control temporal a través de los \emph{sprints}, valorando al inicio de cada uno 
cuál sería la duración adecuada basándonos en las tareas que se iban a realizar y el producto que se esperaba al final de esa iteración. 
\item
Seguimiento de tareas en un panel \emph{Kanban}.
\end{itemize}

Además de los \emph{sprints} se diseñaron hitos relevantes denominados \emph{Milestones} en \emph{GitHub} que sirven como referencia de la producción realizada. 

\subsection{Sprints}

A continuación se muestra un resumen de los \emph{sprints} que han tenido lugar en las diferentes fases de este trabajo:


\subsubsection{Sprint 0. Presentación inicial a tutores del TFG}

\imagen{sprint_00.png}{Sprint 0}{1}

\begin{itemize}
\item  
\textbf{Duración:} 19/10/2023 a 23/10/2023

\item
\textbf{Objetivo:} Presentar una idea de TFG a los posibles tutores. 

\item
\textbf{Contexto:} Inicio del último curso del Grado de Ingeniería Informática. Se trata de contactar con los posibles tutores que controlen el desarrollo de un TFG. Este TFG estará basado en la idea de desarrollar una herramienta, escrita en Python, para realizar análisis financiero. 

\item
\textbf{Tareas:}
	\begin{itemize}
	\tightlist
	\item 
	Contactar con los posibles tutores por email.
	\item
	Contactar por Teams con los tutores que muestren interés en el TFG e intercambiar primeras impresiones.
	\item
	Determinar las herramientas necesarias para desarrollar un TFG.
	\item
	Preguntar por consejos y opiniones que puedan ser relevantes. 
	\item
	Concertar cita para una futura reunión cuando se haya avanzado con los primeros pasos del TFG.
	\end{itemize}
\end{itemize}



\subsubsection{Sprint 1. Empezar TFG con selección de herramientas adecuadas}

\imagen{sprint_01.png}{Sprint 1}{1}

\begin{itemize}
\item  
\textbf{Duración:} 24/10/2023 a 08/11/2023

\item
\textbf{Objetivo:} Empezar el desarrollo del TFG y su correspondiente documentación. Hacerlo utilizando las herramientas adecuadas sugeridas por los tutores en el trascurso del sprint 0. 

\item
\textbf{Contexto:} Tras haber mantenido una primera reunión con los tutores, se detectan ciertas necesidades en cuanto a la utilización de herramientas adecuadas para el desarrollo de un TFG. 

\item
\textbf{Tareas:}
	\begin{itemize}
	\tightlist
	\item 
	Lectura completa de documentación de TFGs disponible en la plataforma UBUVirtual. 
	\item 	
	Rellenar el formulario de \emph{Oferta de TFG Grado de Informática Online}. Buscar referencias bibliográficas teóricas más adecuadas. 
	\item 
	Utilización de \emph{Zube}\citep{online:zube} para gestionar un proyecto con metodología ágil. Para ello, hay que buscar la documentación adecuada que permita entender su utilización.
  	\item 
  	Crear un repositorio en \emph{GitHub} que permita llevar un seguimiento de las acciones realizadas. 
  	\item 
  	En el repositorio, añadir a los tutores e integrar con \emph{Zube}. 
	\item 
	Determinar la mejor herramienta para la documentación del TFG. Elegir entre \LaTeX y \emph{Word}. Justificar la elección y documentar. 
  	\item 
  	Buscar un gestor de referencias bibliográficas: \emph{Zotero vs Mendeley}. 
  	\item 
  	Concertar cita para una futura reunión cuando se haya avanzado con los primeros pasos del TFG.
  	\end{itemize}
\end{itemize}


\subsubsection{Sprint 2. Mostrar primeras tablas y gráficos en servidor web}

\imagen{sprint_02.png}{Sprint 2}{1}

\begin{itemize}
\item  
\textbf{Duración:} 08/11/2023 a 29/11/2023

\item
\textbf{Objetivo:} Comenzar a desarrollar una aplicación web, mostrando información de precios cotizados y de gráficos. Almacenar información en una base de datos SQLite3, no en DataFrames ni en archivos .csv. 

\item
\textbf{Contexto:} Una de las primeras recomendaciones de los tutores fue no utilizar archivos .csv para almacenar la información, por lo tanto, se va a utilizar un sistema de almacenamiento en base de datos. La idea general es que dichos datos se descargan con una API, se almacenan y, posteriormente, se procesan para mostrar la información relevante al usuario en archivos HTML. 

\item
\textbf{Tareas:}
	\begin{itemize}
	\tightlist
	\item 
	Preparar entorno local para utilizar una base de datos SQLite3. Instalar una GUI para la BD (no relevante para el usuario final). 
	\item 	
	Preparar clases (Descargador, Vista, Ticker y/o similares) que permitan automatizar la tarea de almacenar información en una base de datos. 
	\item 
	Mostrar un primer producto - no necesariamente visualmente atractivo - en un servidor web con tablas extraídas de la BD. 
  	\item 
  	Documentar, en el capítulo 4 de la memoria, las herramientas utilizadas. 
  	\item 
  	Concertar cita para una futura reunión con los tutores. 
  	\end{itemize}
\end{itemize}


\subsubsection{Sprint 3. Mejorar lógica de web y proteger enlaces con registro. Mostrar gráficos y dar primeros estilos}

\imagen{sprint_03.png}{Sprint 3}{1}

\begin{itemize}
\item  
\textbf{Duración:} 29/11/2023 a 13/12/2023

\item
\textbf{Objetivo:} Mejorar la lógica de acceso a la web. Proteger los enlaces no públicos mediante registro de usuarios. Mostrar gráficos y no sólo tablas en un estilo visualmente atractivo. 

\item
\textbf{Contexto:} Una vez llevado el proyecto a \emph{producción} (subido a un servidor web) hay que mejorar y asegurar los enlaces para que no sean accesibles en un mal uso de los mismos. Además, es necesario hacer un registro de usuarios que permita mostrar información relevante dependiendo del rol (invitado / registrado). 

\item
\textbf{Tareas:}
	\begin{itemize}
	\tightlist
	\item 
	Crear lógica de control de usuarios registrados.
	\item 	
	Asignar permiso de acceso a enlaces dependiendo del tipo de usuario. 
	\item 
	Crear portada y dar primeros estilos visuales. 
  	\item 
  	Preparar un índice o zona de \emph{breadcrumbs} de la página para mejorar la navegación.
  	\item 
  	Diseñar una vista que permita mostrar gráficos según el valor cotizado escogido. 
  	\end{itemize}
\end{itemize}


\subsubsection{Sprint 4. Manejar imágenes estáticas. Añadir base de datos de DJ30 y buscar método para auto-actualizar las BDs}

\imagen{sprint_04.png}{Sprint 4}{0.9}

\begin{itemize}
\item  
\textbf{Duración:} 14/12/2023 a 28/12/2023

\item
\textbf{Objetivo:} Mejorar estética y mostrar imágenes estáticas. Crear y añadir una nueva base de datos. Hacer las bases de datos auto-actualizables.

\item
\textbf{Contexto:} El proyecto en local permite mostrar imágenes estáticas en la navbar y en portada, pero en \emph{producción} (subido a un servidor web) no. Hay que mejorar este aspecto para que resulte visualmente atractivo. Además, hay que crear y añadir una nueva base de datos para el índice DJ30 y hay que buscar un método para auto-actualizar las BD en producción: con \emph{cron}\citep{wiki:cron} o con \emph{tasks}. 

\item
\textbf{Tareas:}
	\begin{itemize}
	\tightlist
	\item 
	Mostrar imágenes estáticas en producción (en servidor web).
	\item 	
	Crear una nueva BD y hacerla accesible. 
	\item 
	Investigar e implantar un método de auto-actualización de BDs en producción. 
  	\end{itemize}
\end{itemize}


\subsubsection{Sprint 5. Añadir un carrusel de noticias y cambiar el acceso a las BDs. Incorporar tests y mejorar los logs}

\imagen{sprint_05.png}{Sprint 5}{1}

\begin{itemize}
\item  
\textbf{Duración:} 04/01/2024 a 18/01/2024

\item
\textbf{Objetivo:} Añadir noticias relacionadas con los mercados financieros, para hacer un interfaz más útil. Cambiar la lógica de acceso a las bases de datos, para evitar las sentencias en SQL y aprovechar las capacidades del patrón MVT\citep{online:django_MVT_2} de Django. Aplicar primeros tests. 

\item
\textbf{Contexto:} Ahora mismo la portada principal es una imagen estática que permite \emph{login}, pero que no aporta la utilidad. Parece razonable añadir una zona de noticias, que hagan la página más dinámica. Además, aunque tengo diseñados unos pequeños tests, se tiene que mejorar toda la estructura de \emph{testing} y \emph{log}. 

\item
\textbf{Tareas:}
	\begin{itemize}
	\tightlist
	\item 
	Añadir una nueva aplicación al proyecto, que permita controlar lo que ocurre en la portada:
		\begin{itemize}
		\tightlist
		\item
		Mostrar noticas en portada.
		\item
		Mostrar información sobre los mejores y peores stocks.
		\item
		Mejorar el logging para lo que ocurra en esta nueva aplicación.
		\end{itemize}
	\item
	Mejorar la lógica de acceso a las BDs. Aplicar MVC y evitar consultas directas de SQL en la medida de lo posible.
	\item 	
	Ampliar los tests y mejorar la estructura de testing. 
  	\end{itemize}
\end{itemize}


\subsubsection{Sprint 6. Añadir un \emph{dashboard} para almacenar y controlar una cartera por parte del usuario. Operaciones CRUD en BDs}

\imagen{sprint_06.png}{Sprint 6}{0.85}

\begin{itemize}
\item  
\textbf{Duración:} 22/01/2024 a 05/02/2024

\item
\textbf{Objetivo:} Añadir un área personal que permita realizar el control de una cartera de acciones. Poder crear informes con la composición de la cartera y mostrar porcentajes de inversión.   

\item
\textbf{Contexto:} Aunque un usuario puede consultar los datos de un valor concreto y de su índice de referencia, ahora mismo no puede tener una lista de seguimiento. Se pretende integrar una nueva \emph{app} que permita ofrecer un \emph{dashboard} a los usuarios. 

\item
\textbf{Tareas:}
	\begin{itemize}
	\tightlist
	\item 
	Añadir una nueva \emph{app} que permita controlar un \emph{DashBoard}. 
		\begin{itemize}
		\tightlist
		\item
		Permitir almacenar información sobre las inversiones. 
		\item
		Mostrar acciones compradas, i.e., posiciones abiertas.
		\item
		Mostrar una lista de seguimiento. 
		\item
		Mejorar el \emph{logging} para lo que ocurra en esta nueva aplicación.
		\item
		Preparar los tests para el control de esta nueva \emph{app}. 
		\item
		Integrar esta aplicación con las actuales. 
		\end{itemize}
	\item
	Aprovechar la lógica mejorada de acceso a las BDs para facilitar todas las operaciones CRUD. 
	\item 	
	Ampliar los tests.
  	\end{itemize}
\end{itemize}


\subsubsection{Sprint 7. Refactoring y testing. Poder eliminar stocks en seguimiento desde el DashBoard.}

\imagen{sprint_07.png}{Sprint 7}{0.85}

\begin{itemize}
\item  
\textbf{Duración:} 07/02/2024 a 14/02/2024

\item
\textbf{Objetivo:} Ahora mismo hay bastantes casos sin cubrir con los tests, especialmente en los \emph{views.py}. Antes de continuar, hay que realizar un testeo exhaustivo para asegurar una buena base en desarrollos posteriores. Además, voy a incorporar una nueva funcionalidad dentro del \emph{DashBoard}, que permitirá eliminar valores en seguimiento.

\item
\textbf{Contexto:} Tras intentar seguir un enfoque TDD en la creación del DashBoard en el sprint 6, se detecta la necesidad de resolver algunos fallos y cubrir el código, al 100\%, con tests. Además, hay que mejorar el estilo y la calidad del código, así como medir dicha calidad con alguna herramienta que ofrezca información final.

\item
\textbf{Tareas:}
	\begin{itemize}
	\tightlist
	\item 
	Cubrir todas las normas de estilo de \emph{PEP-8} enhancement
	\item
	Solucionar problema formularios \emph{DashBoard}.
	\item
	Incluir funcionalidad de eliminación de stocks en seguimiento.
	\item
	Completar tests unitarios de \texttt{News.views}
	\item
	Completar tests unitarios \texttt{DashBoard.views}
	\item
	Completar tests unitarios de \texttt{Analysis.views}
	\item
	Medir con \emph{flake8} y \emph{pylint} la calidad del código.
  	\end{itemize}
\end{itemize}


\subsubsection{Sprint 8. Correlaciones entre valores con NetworkX}

\imagen{sprint_08.png}{Sprint 8}{0.85}

\begin{itemize}
\item  
\textbf{Duración:} 17/02/2024 a 02/03/2024

\item
\textbf{Objetivo:} Permitir al usuario comparar la evolución de un valor con respecto al resto, en un tiempo determinado. 

La idea es mostrar al usuario un grafo creado con \emph{NetworkX}\citep{online:networkx} que permita entender las correlaciones (positiva y negativa) que haya entre diferentes valores. De esta manera se podrá hacer un análisis de diversificación de cartera, ya que tener valores muy correlacionados puede ser indicativo de no tener la cartera bien balanceada. 

\item
\textbf{Contexto:} El usuario puede acceder a la información de un valor pero no puede compararlo con otros. Además, aunque tenga valores comprados o en seguimiento, realmente, ahora mismo, no se le está ofreciendo capacidad analítica, así que hay que proporcionársela. 


\item
\textbf{Tareas:}
	\begin{itemize}
	\tightlist
	\item 
	Añadir dos índices adicionales, con sus valores, para ampliar la información disponible y ofrecer un mayor abanico comparativo. 
	\item
	Agregar una funcionalidad que permita ver la correlación de un valor con todos los demás disponibles. 
		\begin{itemize}
		\tightlist
		\item 
		Como la correlación puede ser positiva o negativa, hacer diferenciación entre ambas. 
		\item
		Además, puede que interese obtener información sobre las mayores correlaciones (positivas o directas; y negativas o inversas)
		\item
		Una vez obtenidas las correlaciones, sobre los precios de cierre, ver la información en grafos de \emph{NetworkX}. 
		\end{itemize}
	\item
	Adicionalmente, intentar permitir al usuario ver los gráficos de la mayor correlación de forma conjunta. 
	\item
	Ampliar los tests a la nueva funcionalidad. 
	\item
	Integrar los grafos en la plataforma web de forma temporal para cada usuario (sin almacenar en el servidor).
  	\end{itemize}
\end{itemize}

\newpage
\subsubsection{Sprint 9. Aplicación de modelos. \emph{Forecasting} para series temporales}

\imagen{sprint_09.png}{Sprint 9}{1}

\begin{itemize}
\item  
\textbf{Duración:} 13/03/2024 a 27/03/2024

\item
\textbf{Objetivo:} Aportar una visión diferenciadora al análisis de valores cotizados, permitiendo que un usuario pueda aplicar un modelo que le permita tener una idea intuitiva de la futura evolución de la cotización. En cualquier caso, se informará al usuario de que los datos aportados por el modelo no pueden ser considerados como una fuente única para tomar decisiones de inversión. 


\item
\textbf{Contexto:} En las webs que ofrecen información de productos cotizados no se ofrece un apartado de experimentación con modelos de \emph{forecasting}. Puede ser un factor diferenciador para aquellos usuarios que entiendan los riesgos y las ventajas de aplicar estos modelos.  


\item
\textbf{Tareas:}
	\begin{itemize}
	\tightlist
	\item 
	Añadir el análisis a través de un modelo ARIMA. 
	\item
	Añadir el análisis a través de una red LSTM. 
	\item
	Permitir al usuario ver alguna información gráfica que facilite la comprensión de los resultados. 
	\item
	Ampliar los tests de las nuevas funcionalidades. 
	\item
	Integrar las gráficas o diagramas en la plataforma web de forma temporal para cada usuario (sin almacenar en el servidor).
  	\end{itemize}
\end{itemize}


\subsubsection{Sprint 10. Documentación y limpieza de código y avanzar considerablemente en la memoria.}

\imagen{sprint_10.png}{Sprint 10}{1}

\begin{itemize}
\item  
\textbf{Duración:} 10/04/2024 a 24/04/2024

\item
\textbf{Objetivo:} Documentación de código y de la memoria. 


\item
\textbf{Contexto:} Actualmente el proyecto está documentado con notas internas del autor y debe empezarse a documentar utilizando las herramientas y guías oficiales. 


\item
\textbf{Tareas:}
	\begin{itemize}
	\tightlist
	\item 
	Configurar \emph{MikTex}.
	\item
	Descargar y preparar plantilla oficial de Trabajos de Fin de Grado. 
	\item
	Añadir resumen / abstract. 
	\item
	Documentar introducción. 
	\item
	Describir objetivos del proyecto. 
	\item
	Documentar conceptos teóricos. 
	\item
	Explicar las técnicas y herramientas utilizadas. 
	\item
	Indicar aspectos relevantes del desarrollo del proyecto que se hayan dado hasta ahora. 
	\item
	Detallar trabajos relacionados.
	\item
	Añadir referencias bibliográficas en todos los apartados de la memoria. 
	\item
	Medir la calidad del código e introducir las mejoras necesarias. 
  	\end{itemize}
\end{itemize}


\subsubsection{Sprint 11. Gráfica de Markowitz con frontera efciente y ratio de Sharpe. Mejoras visuales.}

\imagen{sprint_11.png}{Sprint 11}{1}

\begin{itemize}
\item  
\textbf{Duración:} 09/05/2024 a 19/05/2024

\item
\textbf{Objetivo:} Favorecer el estudio de una cartera de valores a través de la distribución de pesos de los mismos. Mostrar una gráfica de Markowitz con la cartera ideal encontrada por simulación de Montecarlo y, también, por optimización de funciones. Realizar pequeñas mejoras estéticas y funcionales. 


\item
\textbf{Contexto:} Es necesario mejorar el \emph{DashBoard} de los usuarios para que reciban una información más completa y, de entre las opciones barajadas, una de las más interesantes es la de mostrar cuál sería la mejor distribución de los valores que ya tenga en cartera el usuario, basándome en los retornos y pesos en la propia cartera. 


\item
\textbf{Tareas:}
	\begin{itemize}
	\tightlist
	\item 
	Añadir nuevos métodos que permitan hacer una simulación de Montecarlo con diferentes pesos de valores en una cartera. 
	\item
	Nuevos métodos en \emph{DashBoard} para calcular frontera eficiente. 
	\item
	Desarrollar las funciones de optimización que permitan calcular la mejor distribución de pesos.
	\item
	Preparar métodos para mostrar gráfica de Markowitz junto con la frontera eficiente y las mejores carteras encontradas. 
	\item
	Implementar los tests necesarios para las nuevas funcionalidades. 
	\item
	Integrar en las plantillas existentes del \emph{DashBoard} para mostrar. 
	\item
	Documentar de forma interna en el apartado de conceptos teóricos (apartado 3 de la memoria del TFG). 
	\item
	Mejoras estéticas en botones y página principal. 
	\item
	Añadir referencias bibliográficas en todos los apartados de la memoria aunque sea a nivel local. 
	\item
	Medir la calidad y testear con \emph{GitHub actions}. 
	\item
	Desplegar nueva documentación de código con \emph{Sphinx}. 
  	\end{itemize}
\end{itemize}



\subsubsection{Sprint 12. Mejorar gráficas ARIMA. Incluir estrategias basadas en \emph{machine learning}}

\imagen{sprint_12.png}{Sprint 12}{0.9}

\begin{itemize}
\item  
\textbf{Duración:} 22/05/2024 a 26/05/2024

\item
\textbf{Objetivo:} Completar documentación teórica de ARIMA en memoria, añadir apartado de \emph{trading} algorítmico con su documentación y contemplar redes LSTM como posible mejora del proyecto. 


\item
\textbf{Contexto:} Ahora mismo se puede acceder a un apartado de \emph{forecasting} con redes LSTM pero no me convencen los resultados obtenidos, así que se dejará como una de las posibles mejoras del proyecto. Sin embargo, voy a añadir un apartado nuevo al \emph{Lab} sobre \emph{trading} algorítmico. 

Además, tengo que terminar la documentación de ARIMA en la memoria y añadir la nueva documentación de \emph{trading} algorítmico. 


\item
\textbf{Tareas:}
	\begin{itemize}
	\tightlist
	\item 
	Terminar documentación ARIMA. 
	\item
	Mejorar gráficas y plantillas HTML para ARIMA.  
	\item
	Incluir mejoras de estrategia basada en ML. 
	\item
	Pasar redes LSTM a apartado de posibles mejoras del proyecto. 
	\item
	Realizar nuevos tests para estrategias basadas en ML.
	\item
	Crear un \emph{release} con el cambio, porque es relevante. 
  	\end{itemize}
\end{itemize}



\subsubsection{Sprint 13. Finalizar memoria. Empezar anexos. Correcciones en plantillas y en \emph{cron} de servidor web}

\imagen{sprint_13.png}{Sprint 13}{0.9}

\begin{itemize}
\item  
\textbf{Duración:} 31/05/2024 a 04/06/2024

\item
\textbf{Objetivo:} Finalizar memoria y empezar anexos. Realizar mejoras estéticas. Corregir tarea \emph{cron} en servidor web. Realizar la \emph{release} que está pendiente desde el \emph{sprint} anterior.


\item
\textbf{Contexto:} En este momento es posible trabajar con todas las herramientas, pero quedan algunos detalles que mejorar antes de realizar un \emph{release}. Además, es necesario avanzar con la documentación.


\item
\textbf{Tareas:}
	\begin{itemize}
	\tightlist
	\item 
	Finalizar secciones 5, 6 y 7 de memoria.
	\item
	Empezar anexos. 
	\item
	Mejorar gráficas y plantillas para estrategias basadas en ML. 
	\item
	Ampliar tests para cubrir todos los métodos. 
	\item
	Mejorar la calidad de código. Comprobar con \emph{pylint}.  
	\item
	Mejoras gráficas en grafos de correlaciones de \emph{NetworkX} y adaptación de plantillas en local.
	\item
	Corregir problemas de tarea \emph{cron} en servidor web. 
	\item
	Crear un release con los cambios, siempre y cuando haya terminado todas las correcciones.
  	\end{itemize}
\end{itemize}

\subsubsection{Sprint 14. Finalizar anexos. Grabar vídeos. Realizar pequeñas mejoras}

\imagen{sprint_14.png}{Sprint 14}{0.9}

\begin{itemize}
\item  
\textbf{Duración:} 05/06/2024 a 11/06/2024

\item
\textbf{Objetivo:} Finalizar los anexos, grabar los vídeos y realizar pequeñas mejoras estéticas y de calidad. Realizar la entrega del TFG.. 


\item
\textbf{Contexto:} Se dispone de la última versión del proyecto en \emph{GitHub} y en la web, pero es necesaria la finalización de los anexos y realizar pequeñas mejoras, especialmente en cuanto a calidad de código y documentación de métodos. 


\item
\textbf{Tareas:}
	\begin{itemize}
	\tightlist
	\item 
	Finalizar aptdo. 2 de anexos. 
	\item
	Finalizar aptdo. 3 de anexos. 
	\item
	Finalizar aptdo. 4 de anexos. 
	\item
	Finalizar aptdo. 5 de anexos. 
	\item
	Finalizar aptdo. 6 de anexos. 
	\item
	Grabar vídeos.
	\item
	Mejoras en calidad de código.
	\item
	Mejoras en documentación de métodos.
	\item
	Realizar mejoras estéticas si da tiempo: cambiar presentación de datos de un valor individual y mejorar grafos de NetworkX (tarea no prioritaria). 
	\item
	Crear un \emph{release} final.
  	\end{itemize}
\end{itemize}


\section{Estudio de viabilidad}

\subsection{Viabilidad económica}

Se va a realizar una estimación de costes lo más aproximada a la realidad, como si el proyecto se hubiese realizado en un entorno empresarial. 

\subsubsection{Estudio de costes}

De manera general se pueden considerar costes de personal, de \emph{hardware} y \emph{software} y, en caso de necesidad, costes de infraestructura. La mayor partida económica se tiene que destinar a los recursos humanos:

\tablaSmallSinColores{Costes en RRHH}{l|r|r}{costes_personal}{
\textbf{Concepto} & \textbf{Coste anual} & \textbf{Prorrateo (6 meses)}\\
}{
\textit{Salario bruto} & 25.000,00€ & 12.500,00€\\
\textit{Retención IRPF} & 4.250,00€ & 2.125,00€\\
\textit{Seguridad Social} & 1.587,50€ & 793,75€\\
\textit{Salario neto} & 19.162,50€ & 9.581,25€\\ 
}

Se aplica un 17\% de retención sobre la nómina, considerando ausencia de deducciones tributarias.

A continuación se muestran los costes de \emph{hardware} y \emph{software}. Para el \emph{hardware} se estima una amortización de 5 años y una utilización de 6 meses. El \emph{software}, por su parte, tendrá una amortización estimada de 2 años. 

\tablaSmallSinColores{Costes de \emph{HW} y \emph{SW}}{l|r|r}{costes_hw_sw}{
\textbf{Concepto} & \textbf{Coste} & \textbf{Coste amortizado} \\
}{
\textit{Ordenador portátil} & 1.600,00€ & 160,00€\\
\textit{Licencia MS Windows 10 pro} & 279,00€ & 69,75€\\
}

Por último, se realiza un análisis de costes de infraestructura (solo asimilable en caso de necesidad) de nuevo, el coste amortizado se calcula a 6 meses:

\tablaSmallSinColores{Costes infraestructura}{l|r|r}{costes_infraestructura}{
\textbf{Concepto} & \textbf{Coste anual} & \textbf{Coste amortizado} \\
}{
\textit{Consumo eléctrico} & 350,00€ & 175,00€\\
\textit{Espacio \emph{coworking}} & 4.000,00€ & 2.000,0€\\
\textit{Material de oficina} & 10,00€ & 5,00€\\
\textit{Alojamiento web} & 72,00€ & 36,00€\\
}


Los costes totales estimados son:

\tablaSmallSinColores{Costes totales}{l|r}{coste_total_estimado}{
\textbf{Concepto} & \textbf{Coste}\\
}{
\textit{Costes en RRHH} & 12.500,00€\\
\textit{Costes de \emph{HW} y \emph{SW}} & 229.75€\\
\textit{Costes infraestructura} & 2.216,00€\\
\textit{\textbf{TOTAL}} & 14.945,75€\\
}

\subsubsection{Análisis de beneficios}

Este trabajo no se plantea inicialmente como un mecanismo para generar beneficios económicos, sino como una herramienta gratuita para ayudar a los inversores a mejorar la distribución de sus carteras y para que los analistas financieros experimenten con diferentes técnicas. 

En caso de monetizar el proyecto se podría optar por incluir publicidad en la web para generar ingresos que, al menos, cubran los gastos anuales. 

\subsection{Viabilidad legal}

\subsubsection{Descargo de responsabilidades}

El \emph{software} que se proporciona en este trabajo es solo para fines informativos y no debe considerarse como asesoramiento financiero. La información proporcionada no está garantizada como precisa o completa y puede cambiar sin previo aviso.

El usuario es el único responsable de sus decisiones de inversión y de las consecuencias de las mismas. No se debe confiar en este \emph{software} para tomar decisiones de inversión sin realizar una investigación profunda y consultar con un asesor financiero calificado.

\subsubsection{Licencias \emph{software}}

En este trabajo se utilizan múltiples librerías de terceros. Además, aunque los datos son almacenados en bases de datos propias, hay que considerar que se hace uso de información externa a través de la API \emph{yFinance}. 

Las licencias asociadas a las librerías y APIs utilizadas, se detallan en la siguiente tabla:

\tablaSmallSinColores{Licencias de terceros}{l|r}{licencias_terceros}{
\textbf{Librería / API} & \textbf{Licencia}\\
}{
\textit{Python} & OSI-Open Source\\
\textit{Django} & BSD\\
\textit{yFinance} & Apache 2.0\\
\textit{Pandas} & BSD\\
\textit{Numpy} & BSD 3-Clause\\
\textit{Scikit-Learn} & BSD 3-Clause\\
\textit{Scipy} & BSD\\
\textit{NetworkX} & BSD 3-Clause\\
\textit{Matplotlib} & Licencia libre propia\\
\textit{Plotly} & MIT\\
\textit{Statsmodels} & BSD 3-Clause\\
\textit{Feedparser} & BSD 3-Clause\\
\textit{NewsAPI} & MIT\\
\textit{pmdarima} & MIT\\
\textit{Keras} & Apache 2.0\\
}

Todas las librerías y el \emph{framework Django} utilizan licencias de código abierto permisivas - la más restrictiva sería Apache 2.0 - lo que significa que el código de \emph{FAT: Financial Analysis Tool} se puede usar, modificar y distribuir libremente, incluyendo para fines comerciales, siempre que se cumplan los requisitos de atribución y exención de responsabilidad establecidos en cada licencia.

Por tanto, se toma la decisión de publicar este proyecto bajo licencia CC BY-NC-SA 4.0\citep{online:licencia}.

\subsubsection{Consideraciones adicionales}

Dada la condición académica de este proyecto no hay limitación en el uso de los datos financieros almacenados y no se recoge información personal; pero dependiendo del uso que se dé a este \emph{software} es posible que se deba cumplir con regulaciones específicas, como la Ley de Protección de Datos de la Unión Europea (GDPR) o la Ley de Protección de la Privacidad de la Información Financiera (FINRA). 













\apendice{Especificación de Requisitos}

\section{Introducción}

En este anexo se recogen los requisitos y casos de uso que se han tenido en cuenta para el desarrollo de este trabajo. Se han seguido las recomendaciones del estándar IEEE 830-1998\citep{1998IEEESpecifications} como guía de buenas prácticas.

En este proyecto se ha tratado de crear una especificación de requisitos que cumpla con las siguientes condiciones:

\begin{itemize}
\tightlist
\item
	Completa y consistente
	\begin{itemize}
	\item 
	Incluir todos los requisitos y referencias necesarias.
	\item
	Ser coherente con los propios requisitos y otros documentos de especificación.
	\end{itemize}
\item
	Claridad y accesibilidad
	\begin{itemize}
	\item 
	Redacción clara para evitar malas interpretaciones.
	\item
	Uso de términos y definiciones precisos.
	\end{itemize}
\item
	Verificabilidad y trazabilidad
	\begin{itemize}
	\item 
	Debe de existir un método finito y sin costo para probar los requisitos.
	\item
	Uso de términos y definiciones precisos.
	\end{itemize}
\item
	Modificabilidad y priorización
	\begin{itemize}
	\item 
	Fácilmente modificable.
	\item
	Jerarquía de priorización según relevancia para el negocio (o fin).
	\end{itemize}
\item
	Correctitud y accesibilidad
	\begin{itemize}
	\item 
	El software debe cumplir con los requisitos de la especificación.
	\item
	Accesibilidad y facilidad de comprensión para los usuarios y desarrolladores.
	\end{itemize}
\end{itemize}


\section{Objetivos generales}\label{objetivos-generales}

Este trabajo persigue los siguientes objetivos generales:

\begin{itemize}
\item
Proporcionar una herramienta que otorgue capacidad crítica en inversiones personales.
\item
El sistema tiene que poder controlar la evolución de valores en cartera y en seguimiento. 
\item
Tiene que mostrarse información relevante en cuanto a distribución de carteras y posibles pesos para los valores escogidos. 
\item
Permitir trabajar, de manera experimental, con herramientas de \emph{trading} algorítmico. 
\item
Facilitar el acceso a la información con una página web pública y gratuita. 
\end{itemize}


\section{Catálogo de requisitos}

\subsection{Requisitos funcionales}

\begin{itemize}
\item 
\textbf{RF-1 Mostrar portada con información general de diferentes índices bursátiles:} la web tiene que tener información agregada de todos los índices disponibles. 
    \begin{itemize}
        \item \textbf{RF-1.1 Mostrar carrusel de noticias generales}: en la portada se 
        tiene que disponer de noticias relevantes dentro del mundo bursátil.
        \item \textbf{RF-1.2 Mostrar mejores y peores valores de cada índice}: se tiene que 
        poder consultar, de forma ágil, cuáles han sido los mejores y peores valores de 
        cada índice.
        \begin{itemize}
            \item \textbf{RF-1.2.1 Mostrar gráfica}: tiene que estar disponible una gráfica 
            evolutiva junto a cada valor de los mostrados. 
        \end{itemize}
	\end{itemize}

\item
\textbf{RF-2 Control de usuarios}: los usuarios tendrán acceso a información adicional si
están dados de alta. 
	\begin{itemize}
        \item \textbf{RF-2.1 Permitir registro}: el usuario se podrá registrar. 
        \item \textbf{RF-2.2 Permitir hacer \emph{login}}: el usuario podrá acceder con su 
        cuenta en la web.
        \item \textbf{RF-2.3 Permitir \emph{logout}}: el usuario podrá cerrar su sesión. 
	\end{itemize}

\item
\textbf{RF-3 Gestionar \emph{DashBoard}}: tiene que haber una zona de usuario, denominada \emph{DashBoard}, con información relevante sobre valores seleccionados por el usuario. 
	\begin{itemize}
        \item \textbf{RF-3.1 Crear nuevo valor en cartera}: el usuario podrá añadir un 
        nuevo valor, con fecha y precio de compra, a su cartera.
        \begin{itemize}
            \item \textbf{RF-3.1.1 Mostrar 'donut', sectores y divisas}: si un usuario 
            tiene valores en cartera, éste contará con diagramas de información relevante.
            \item \textbf{RF-3.1.2 Ver gráfica de Markowitz, ratio de Sharpe y pesos}: si 
            un usuario tiene valores en cartera podrá ver lla mejor distribución de pesos 
            para la misma. 
            \item \textbf{RF-3.1.3 Ver posiciones abiertas y evolución de cartera}: si hay 
            valores en cartera se podrá comprobar la evolución de la misma. 
        \end{itemize}
		\item \textbf{RF-3.2 Eliminar valor de cartera}: el usuario podrá eliminar un 
        valor de los que tuviera almacenados. 
        \item \textbf{RF-3.3 Crear un nuevo valor en seguimiento}: el usuario podrá 		
        hacer seguimiento de valores de su interés.
        \item \textbf{RF-3.4 Eliminar valor de seguimiento}: el usuario podrá eliminar un 
        valor de los que tuviera en seguimiento. 
	\end{itemize}

\item
\textbf{RF-4 Consultar índice}: se debe presentar, de forma ordenada, información relevante de cada índice por separado.  
	\begin{itemize}
        \item \textbf{RF-4.1 Mostrar tabla de valores}: al consultar un índice se mostrará 
        una tabla con todos los componentes de ese índice, con información relevante sobre 
        la última sesión disponible. 
        \item \textbf{RF-4.2 Consultar un valor del índice}: se podrá consultar información 
        de un único valor. 
        \begin{itemize}
            \item \textbf{RF-4.2.1 Dar acceso a gráfica interactiva}: se ofrecerá una 
            gráfica interactiva, con medias móviles y volumen. 
            \item \textbf{RF-4.2.2 Mostrar distribución de retornos}: habrá una gráfica con 
            la distribución que hayan seguido los retornos en los últimos meses. 
            \item \textbf{RF-4.2.3 Mostrar datos del último mes}: disponer datos en formato 
            tabular para ver precios de apertura y cierre de los últimos 30 días. 
            \item \textbf{RF-4.2.4 Ver evolución del sector}: facilitar una comparativa con 
            el sector de referencia del valor. 
            \item \textbf{RF-4.2.5 Comparar con otros valores}: existirá la posibilidad de 
            comparar gráficas de precios de cierre con otros valores. 
            \item \textbf{RF-4.2.6 Mostrar grafos de correlación}: mostrar grafos de alta 
            correlación positiva y negativa con otros valores. 
        \end{itemize}
        \item \textbf{RF-4.3 Mostrar noticias relacionadas}: facilitar enlaces de fuentes 
        RSS que estén relacionadas con el índice.
        \item \textbf{RF-4.4 Ver gráfica de evolución}: mostrar gráfica de evolución del
        índice en su conjunto. 
	\end{itemize}	



\item
\textbf{RF-5 Gestionar \emph{Lab}}: dar acceso a un laboratorio virtual. 
	\begin{itemize}
        \item \textbf{RF-5.1 Trabajar con ARIMA}: permitir realizar estimaciones y búsqueda 
        de parámetros con ARIMA. 
        \begin{itemize}
            \item \textbf{RF-5.1.1 Buscar parámetros (p,d,q) con fnuciones ACF y PACF}: se 
            mostrarán gráficas de funciones para interpretación del usuario. 
            \item \textbf{RF-5.1.2 Introducir (p,d,q) de forma manual}: el usuario podrá 
            introducir los parámetros deseados para aplicar un modelo ARIMA. 
            \item \textbf{RF-5.1.3 Buscar (p,d,q) de forma automática}: se proporcionará 
            una funcionalidad de búsqueda automática de los parámetros más adecuados a un 
            valor.
            \item \textbf{RF-5.1.4 Hacer búsqueda por rejilla para (p,d,q)}: se podrá 
            realizar una búsqueda de parámetros (p,d,q) de entre una serie de posibles 
            valores preestablecidos.             
        \end{itemize}
        \item \textbf{RF-5.2 Trabajar con \emph{trading} algorítmico}: facilitar el acceso 
        a herramientas de \emph{trading} algorítmico.
        \begin{itemize}
            \item \textbf{RF-5.2.1 Usar algoritmo de cruce de medias}: se mostrará el 
            resultado de buscar las mejores medias móviles para un valor, en un período de 
            tiempo concreto. 
            \item \textbf{RF-5.2.2 Utilizar estrategia basada en \emph{machine learning}}: 
            permitir interactuar con modelos de regresión y clasificación para estimar la 
            tendencia de la próxima sesión de un índice.
        \end{itemize}
	\end{itemize}

\end{itemize}


\subsection{Requisitos no funcionales}

En este apartado se tratará de dar detalle de aquellas características que no son funcionales pero que aportan un valor añadido al proyecto y a la interacción con la herramienta desarrollada:

\begin{itemize}
\item
\textbf{RNF-1 Escalabilidad}: la web tiene que permitir y favorecer, en la medida de lo posible, la incorporación de nuevas funcionalidades y de bases de datos adicionales. 

\item
\textbf{RNF-2 Privacidad}: los datos de usuario, como nombres o contraseñas, se debe gestionar de forma segura. 

\item
\textbf{RNF-3 Disponibilidad}: la web debe ser compatible con los navegadores más modernos y tendrá alta disponibilidad a través de un servicio de alojamiento fiable. 

\item
\textbf{RNF-4 Usabilidad}: la interfaz de la web será \emph{user friendly}, resultar intuitiva y facilitará comentarios de ayuda adicional en los apartados más especializados. 

\item
\textbf{RNF-5 Mantenibilidad}: se tiene que favorecer un mantenimiento posterior a la puesta en producción. Además, se facilitará información a los desarrolladores para que puedan realizar mejoras incrementales posteriores. 
\end{itemize}


\section{Especificación de requisitos}

\subsection{Casos de uso}\label{casos_uso}

A continuación se muestran todos los casos de uso contemplados. Muchos de estos casos de uso se plantearon en las fases iniciales del proyecto y se fueron mejorando en las sucesivas iteraciones de cada \emph{sprint}. Se trató de realizar mejoras incrementales para satisfacer las expectativas de un inversor o potencial cliente:

\imagenSinMargen{img_anex_01_casos_uso.png}{Diagrama de casos de uso. Fuente: elaboración propia}


\begin{table}[p]
	\centering
	\begin{tabularx}{\linewidth}{ p{0.21\columnwidth} p{0.71\columnwidth} }
		\toprule
		\textbf{CU-1}    & \textbf{Consultar portada}\\
		\toprule
		\textbf{Versión}              & 1.0    \\
		\textbf{Autor}                & Rodrigo Merino Tovar \\
		\textbf{Requisitos asociados} & RF-1, RF-1.1, RF-1.2 y RF-1.2.1 \\
		\textbf{Descripción}          & Permite al usuario ver la página principal\\
		\textbf{Precondición}         & La base de datos se encuentra disponible \\
		\textbf{Acciones}             &
		\begin{enumerate}
			\def\labelenumi{\arabic{enumi}.}
			\tightlist
			\item El usuario accede a la página web.
			\item Se muestran noticias del mundo bursátil. 
			\item Se listan los índices con los mejores y peores valores
		\end{enumerate}\\
		\textbf{Postcondición}        & El número de valores mostrados es múltiplo del 
										número de índices disponibles.  \\
		\textbf{Excepciones}          & 
		\begin{itemize}
			\tightlist
			\item Error al cargar valores (mensaje de error).
			\item Error al cargar una noticia (pasar a siguiente). 
		\end{itemize}\\
		\textbf{Importancia}          & Alta \\
		\bottomrule
	\end{tabularx}
	\caption{CU-1 Consultar portada.}
\end{table}


\begin{table}[p]
	\centering
	\begin{tabularx}{\linewidth}{ p{0.21\columnwidth} p{0.71\columnwidth} }
		\toprule
		\textbf{CU-2}    & \textbf{Mostrar carrusel noticias generales}\\
		\toprule
		\textbf{Versión}              & 1.0    \\
		\textbf{Autor}                & Rodrigo Merino Tovar \\
		\textbf{Requisitos asociados} & RF-1, RF-1.1 \\
		\textbf{Descripción}          & Mostrar carrusel de noticias bursátiles\\
		\textbf{Precondición}         & La API \emph{NewsAPI} está activa \\
		\textbf{Acciones}             &
		\begin{enumerate}
			\def\labelenumi{\arabic{enumi}.}
			\tightlist
			\item El usuario ve texto preliminar de una noticia y su imagen asociada. 
			\item Se \emph{navega} entre noticias con selectores laterales o se dejan pasar automáticamente.
			\item Se muestra un botón para \emph{leer más} en la fuente original.
		\end{enumerate}\\
		\textbf{Postcondición}        & La imagen y el enlace externo se corresponden con el texto de la noticia.  \\
		\textbf{Excepciones}          & 
		\begin{itemize}
			\tightlist
			\item Error al cargar una noticia (pasar a siguiente). 
		\end{itemize}\\
		\textbf{Importancia}          & Media \\
		\bottomrule
	\end{tabularx}
	\caption{CU-2 Mostrar carrusel noticias generales.}
\end{table}


\begin{table}[p]
	\centering
	\begin{tabularx}{\linewidth}{ p{0.21\columnwidth} p{0.71\columnwidth} }
		\toprule
		\textbf{CU-3}    & \textbf{Mostrar mejores y peores valores}\\
		\toprule
		\textbf{Versión}              & 1.0    \\
		\textbf{Autor}                & Rodrigo Merino Tovar \\
		\textbf{Requisitos asociados} & RF-1, RF-1.2 y RF-1.2.1 \\
		\textbf{Descripción}          & Mostrar 3 mejores y 3 peores valores de cada índice en la última sesión bursátil.\\
		\textbf{Precondición}         & Las bases de datos están disponibles. \\
		\textbf{Acciones}             &
		\begin{enumerate}
			\def\labelenumi{\arabic{enumi}.}
			\tightlist
			\item El usuario ve valores de cierre de última sesión de mejores valores de un 			índice. 
			\item El usuario ve valores de cierre de última sesión de peores valores de un 				índice. 
			\item Se muestra gráfica asociada a cada valor.
		\end{enumerate}\\
		\textbf{Postcondición}        & Las gráficas se corresponden con los valores 												mejores y peores y se muestran en orden de mejor a 											peor  \\
		\textbf{Excepciones}          & 
		\begin{itemize}
			\tightlist
			\item Error al cargar un valor (mensaje de error). 
		\end{itemize}\\
		\textbf{Importancia}          & Media \\
		\bottomrule
	\end{tabularx}
	\caption{CU-3 Mostrar mejores y peores valores.}
\end{table}


\begin{table}[p]
	\centering
	\begin{tabularx}{\linewidth}{ p{0.21\columnwidth} p{0.71\columnwidth} }
		\toprule
		\textbf{CU-4}    & \textbf{Mostrar gráfica}\\
		\toprule
		\textbf{Versión}              & 1.0    \\
		\textbf{Autor}                & Rodrigo Merino Tovar \\
		\textbf{Requisitos asociados} & RF-1, RF-1.2, RF-1.2.1 y RF-4.4 \\
		\textbf{Descripción}          & Mostrar gráficas de valores o índices.\\
		\textbf{Precondición}         & Las bases de datos están disponibles. \\
		\textbf{Acciones}             &
		\begin{enumerate}
			\def\labelenumi{\arabic{enumi}.}
			\tightlist
			\item El usuario ve gráficas de última sesión de mejores valores de un 						índice. 
			\item El usuario ve gráficas de última sesión de peores valores de un 						índice. 
			\item Se muestran valores de cierre asociados a cada valor.
			\item El usuario consulta gráfica de índice.
		\end{enumerate}\\
		\textbf{Postcondición}        & Las gráficas se corresponden con los valores 												mejores y peores y se muestran en orden de mejor a peor. O la gráfica del índice consultado está disponible. \\
		\textbf{Excepciones}          & 
		\begin{itemize}
			\tightlist
			\item Error al cargar un valor o índice (mensaje de error). 
		\end{itemize}\\
		\textbf{Importancia}          & Media \\
		\bottomrule
	\end{tabularx}
	\caption{CU-4 Mostrar gráfica.}
\end{table}


\begin{table}[p]
	\centering
	\begin{tabularx}{\linewidth}{ p{0.21\columnwidth} p{0.71\columnwidth} }
		\toprule
		\textbf{CU-5}    & \textbf{Gestionar usuario}\\
		\toprule
		\textbf{Versión}              & 1.0    \\
		\textbf{Autor}                & Rodrigo Merino Tovar \\
		\textbf{Requisitos asociados} & RF-2, RF-2.1, RF-2.2 y RF-2.3 \\
		\textbf{Descripción}          & Controlar sesión de usuario.\\
		\textbf{Precondición}         & La base de datos de usuarios está disponible. \\
		\textbf{Acciones}             &
		\begin{enumerate}
			\def\labelenumi{\arabic{enumi}.}
			\tightlist
			\item El usuario se registra. 
			\item El usuario hace \emph{login}. 
			\item Diferentes acciones dependiendo de la labor a realizar.
			\item El usuario hace \emph{logout} para cerrar sesión.
		\end{enumerate}\\
		\textbf{Postcondición}        & Tras \emph{logout} la cookie de sesión se inhabilita. \\
		\textbf{Excepciones}          & 
		\begin{itemize}
			\tightlist
			\item Error al registrar por contraseña o nombre (mensaje informativo). 
			\item Error en \emph{login} por contraseña o nombre (mensaje informativo). 
		\end{itemize}\\
		\textbf{Importancia}          & Alta \\
		\bottomrule
	\end{tabularx}
	\caption{CU-5 Gestionar usuario.}
\end{table}


\begin{table}[p]
	\centering
	\begin{tabularx}{\linewidth}{ p{0.21\columnwidth} p{0.71\columnwidth} }
		\toprule
		\textbf{CU-6}    & \textbf{Registro}\\
		\toprule
		\textbf{Versión}              & 1.0    \\
		\textbf{Autor}                & Rodrigo Merino Tovar \\
		\textbf{Requisitos asociados} & RF-2 y RF-2.1 \\
		\textbf{Descripción}          & Permitir registro de usuarios.\\
		\textbf{Precondición}         & El usuario no está registrado. \\
		\textbf{Acciones}             &
		\begin{enumerate}
			\def\labelenumi{\arabic{enumi}.}
			\tightlist
			\item El usuario se registra. 
			\item Mostrar un mensaje de bienvenida. 
			\item Dejar usuario ya activo con \emph{login}.
			\item Redirigir a página principal.
		\end{enumerate}\\
		\textbf{Postcondición}        & Nuevo usuario en base de datos. \\
		\textbf{Excepciones}          & 
		\begin{itemize}
			\tightlist
			\item Error al registrar por contraseña o nombre (mensaje informativo). 
			\item Error por usuario duplicado (mensaje informativo). 
		\end{itemize}\\
		\textbf{Importancia}          & Alta \\
		\bottomrule
	\end{tabularx}
	\caption{CU-6 Registro.}
\end{table}


\begin{table}[p]
	\centering
	\begin{tabularx}{\linewidth}{ p{0.21\columnwidth} p{0.71\columnwidth} }
		\toprule
		\textbf{CU-7}    & \textbf{Login}\\
		\toprule
		\textbf{Versión}              & 1.0    \\
		\textbf{Autor}                & Rodrigo Merino Tovar \\
		\textbf{Requisitos asociados} & RF-2 y RF-2.2 \\
		\textbf{Descripción}          & Permitir \emph{login} de usuarios.\\
		\textbf{Precondición}         & El usuario no está logueado. \\
		\textbf{Acciones}             &
		\begin{enumerate}
			\def\labelenumi{\arabic{enumi}.}
			\tightlist
			\item El usuario hace \emph{login}. 
			\item Mostrar \emph{DashBoard}. 
			\item Permitir acceso a funciones adicionales.
		\end{enumerate}\\
		\textbf{Postcondición}        & Usuario logueado y con acceso a funciones adicionales. \\
		\textbf{Excepciones}          & 
		\begin{itemize}
			\tightlist
			\item Error al loguear (mensaje informativo y repetición). 
			\item Error forzado (mensaje informativo). 
		\end{itemize}\\
		\textbf{Importancia}          & Alta \\
		\bottomrule
	\end{tabularx}
	\caption{CU-6 Login.}
\end{table}


\begin{table}[p]
	\centering
	\begin{tabularx}{\linewidth}{ p{0.21\columnwidth} p{0.71\columnwidth} }
		\toprule
		\textbf{CU-8}    & \textbf{Logout}\\
		\toprule
		\textbf{Versión}              & 1.0    \\
		\textbf{Autor}                & Rodrigo Merino Tovar \\
		\textbf{Requisitos asociados} & RF-2 y RF-2.3 \\
		\textbf{Descripción}          & Permitir \emph{logout} de usuarios.\\
		\textbf{Precondición}         & El usuario está logueado. \\
		\textbf{Acciones}             &
		\begin{enumerate}
			\def\labelenumi{\arabic{enumi}.}
			\tightlist
			\item El usuario hace \emph{logout}. 
			\item Mostrar página principal. 
			\item No permitir acceso a funciones adicionales.
		\end{enumerate}\\
		\textbf{Postcondición}        & Usuario no logueado y sin acceso a funciones adicionales. \\
		\textbf{Excepciones}          & \\
		\textbf{Importancia}          & Baja \\
		\bottomrule
	\end{tabularx}
	\caption{CU-8 Logout.}
\end{table}


\begin{table}[p]
	\centering
	\begin{tabularx}{\linewidth}{ p{0.21\columnwidth} p{0.71\columnwidth} }
		\toprule
		\textbf{CU-9}    & \textbf{Gestionar \emph{DashBoard}}\\
		\toprule
		\textbf{Versión}              & 1.0    \\
		\textbf{Autor}                & Rodrigo Merino Tovar \\
		\textbf{Requisitos asociados} & RF-3, RF-3.1, RF-3.1.1, RF-3.1.2, RF-3.1.3, RF-3.2, RF-3.3 y RF-3.4  \\
		\textbf{Descripción}          & Mostrar información agregada de valores en cartera y en seguimiento de un usuario.\\
		\textbf{Precondición}         & El usuario está logueado y las bases de datos están disponibles.  \\
		\textbf{Acciones}             &
		\begin{enumerate}
			\def\labelenumi{\arabic{enumi}.}
			\tightlist
			\item El usuario hace \emph{login}. 
			\item Mostrar página de \emph{DashBoard} con información de usuario.
		\end{enumerate}\\
		\textbf{Postcondición}        & El usuario puede consultar precios, añadir o quitar valores en cartera y añadir o quitar valores en seguimiento. \\
		\textbf{Excepciones}          & 
		\begin{itemize}
			\tightlist
			\item No hay valores previos en cartera (se muestran tablas vacías).
			\item No hay valores previos en seguimiento (se muestran tablas vacías).
		\end{itemize} \\
		\textbf{Importancia}          & Alta \\
		\bottomrule
	\end{tabularx}
	\caption{CU-9 \emph{DashBoard}.}
\end{table}


\begin{table}[p]
	\centering
	\begin{tabularx}{\linewidth}{ p{0.21\columnwidth} p{0.71\columnwidth} }
		\toprule
		\textbf{CU-10}    & \textbf{Crear nuevo valor en cartera}\\
		\toprule
		\textbf{Versión}              & 1.0    \\
		\textbf{Autor}                & Rodrigo Merino Tovar \\
		\textbf{Requisitos asociados} & RF-3, RF-3.1, RF-3.1.1, RF-3.1.2 y RF-3.1.3  \\
		\textbf{Descripción}          & Añadir un valor a la cartera del usuario.\\
		\textbf{Precondición}         & El usuario está logueado y las bases de datos están disponibles.  \\
		\textbf{Acciones}             &
		\begin{enumerate}
			\def\labelenumi{\arabic{enumi}.}
			\tightlist
			\item El usuario accede al \emph{DashBoard}. 
			\item Pinchar en botón \emph{Nueva posición de cartera}.
		\end{enumerate}\\
		\textbf{Postcondición}        & El usuario ve un nuevo valor asociado a su cuenta. \\
		\textbf{Excepciones}          & 
		\begin{itemize}
			\tightlist
			\item Valor no existe (mensaje informativo y reintento).
			\item Fecha no válida (mensaje informativo y reintento).
			\item Precio valor no adecuado para la fecha (mensaje informativo y reintento).
			\item Valores formulario fuera de rango (mensaje informativo y reintento).
		\end{itemize} \\
		\textbf{Importancia}          & Alta \\
		\bottomrule
	\end{tabularx}
	\caption{CU-10 Crear nuevo valor en cartera.}
\end{table}


\begin{table}[p]
	\centering
	\begin{tabularx}{\linewidth}{ p{0.21\columnwidth} p{0.71\columnwidth} }
		\toprule
		\textbf{CU-11}    & \textbf{Mostrar \emph{donut}, sectores y divisas}\\
		\toprule
		\textbf{Versión}              & 1.0    \\
		\textbf{Autor}                & Rodrigo Merino Tovar \\
		\textbf{Requisitos asociados} & RF-3, RF-3.1 y RF-3.1.1 \\
		\textbf{Descripción}          & Mostrar información de composición de cartera en gráficas y diagramas que resulten agradables para el usuario.\\
		\textbf{Precondición}         & El usuario está logueado, las bases de datos están disponibles y el usuario tiene valores guardados en cartera.  \\
		\textbf{Acciones}             &
		\begin{enumerate}
			\def\labelenumi{\arabic{enumi}.}
			\tightlist
			\item El usuario accede al \emph{DashBoard}. 
			\item Mostrar \emph{donut} con valores en cartera.
			\item Mostrar diagrama de sectores en los que se está invertido.  
			\item Mostrar inversión por tipo de divisa. 
			\item Ver cambios de divisa de última sesión disponible.
		\end{enumerate}\\
		\textbf{Postcondición}        & El usuario ve información relevante sobre el estado de su cartera. \\
		\textbf{Excepciones}          & 
		\begin{itemize}
			\tightlist
			\item No hay valores previos en cartera (se muestran tablas vacías).
		\end{itemize} \\
		\textbf{Importancia}          & Media \\
		\bottomrule
	\end{tabularx}
	\caption{CU-11 Mostrar \emph{donut}, sectores y divisas.}
\end{table}


\begin{table}[p]
	\centering
	\begin{tabularx}{\linewidth}{ p{0.21\columnwidth} p{0.71\columnwidth} }
		\toprule
		\textbf{CU-12}    & \textbf{Ver gráfica de Markowitz, ratio de Sharpe
y pesos}\\
		\toprule
		\textbf{Versión}              & 1.0    \\
		\textbf{Autor}                & Rodrigo Merino Tovar \\
		\textbf{Requisitos asociados} & RF-3, RF-3.1 y RF-3.1.2 \\
		\textbf{Descripción}          & Mostrar información de distribución de pesos de los valores en cartera.\\
		\textbf{Precondición}         & El usuario está logueado, las bases de datos están disponibles y el usuario tiene valores guardados en cartera.  \\
		\textbf{Acciones}             &
		\begin{enumerate}
			\def\labelenumi{\arabic{enumi}.}
			\tightlist
			\item El usuario accede al \emph{DashBoard}. 
			\item Mostrar gráfica de Markowitz.
			\item Ver simulación de Montecarlo de 10.000 carteras.  
			\item Indicar frontera eficiente. 
			\item Mostrar ratio de Sharpe. 
			\item Mostrar información de distribución de pesos en tablas.
		\end{enumerate}\\
		\textbf{Postcondición}        & El usuario ve información relevante sobre posibles distribuciones mejores para su cartera. \\
		\textbf{Excepciones}          & 
		\begin{itemize}
			\tightlist
			\item No hay valores previos en cartera (se muestran tablas vacías).
		\end{itemize} \\
		\textbf{Importancia}          & Alta \\
		\bottomrule
	\end{tabularx}
	\caption{CU-12 Ver gráfica de Markowitz, ratio de Sharpe y pesos.}
\end{table}


\begin{table}[p]
	\centering
	\begin{tabularx}{\linewidth}{ p{0.21\columnwidth} p{0.71\columnwidth} }
		\toprule
		\textbf{CU-13}    & \textbf{Ver posiciones abiertas y evolución de cartera}\\
		\toprule
		\textbf{Versión}              & 1.0    \\
		\textbf{Autor}                & Rodrigo Merino Tovar \\
		\textbf{Requisitos asociados} & RF-3, RF-3.1 y RF-3.1.3 \\
		\textbf{Descripción}          & Mostrar información de evolución de inversiones realizadas.\\
		\textbf{Precondición}         & El usuario está logueado, las bases de datos están disponibles y el usuario tiene valores guardados en cartera.  \\
		\textbf{Acciones}             &
		\begin{enumerate}
			\def\labelenumi{\arabic{enumi}.}
			\tightlist
			\item El usuario accede al \emph{DashBoard}. 
			\item Mostrar evolución individual de valores en cartera.
			\item Hacer cálculos de cambio de divisa a euros.  
			\item Mostrar evolución de cartera en euros.
		\end{enumerate}\\
		\textbf{Postcondición}        & El usuario ve la evolución de su cartera. \\
		\textbf{Excepciones}          & 
		\begin{itemize}
			\tightlist
			\item No hay valores previos en cartera (se muestran tablas vacías).
		\end{itemize} \\
		\textbf{Importancia}          & Alta \\
		\bottomrule
	\end{tabularx}
	\caption{CU-13 Ver posiciones abiertas y evolución de cartera.}
\end{table}


\begin{table}[p]
	\centering
	\begin{tabularx}{\linewidth}{ p{0.21\columnwidth} p{0.71\columnwidth} }
		\toprule
		\textbf{CU-14}    & \textbf{Eliminar valor de cartera}\\
		\toprule
		\textbf{Versión}              & 1.0    \\
		\textbf{Autor}                & Rodrigo Merino Tovar \\
		\textbf{Requisitos asociados} & RF-3 y RF-3.2 \\
		\textbf{Descripción}          & Eliminar valor guardado en cartera.\\
		\textbf{Precondición}         & El usuario está logueado, las bases de datos están disponibles y el usuario tiene valores guardados en cartera.  \\
		\textbf{Acciones}             &
		\begin{enumerate}
			\def\labelenumi{\arabic{enumi}.}
			\tightlist
			\item El usuario accede al \emph{DashBoard}. 
			\item Mostrar información relativa a cartera.
			\item El usuario pulsa el botón \emph{Eliminar posición de cartera}.  
			\item Mostrar lista de valores en cartera.
			\item Seleccionar valores a eliminar.
			\item Pulsar \emph{Eliminar seleccionados}.
			\item Volver al \emph{DashBoard}
		\end{enumerate}\\
		\textbf{Postcondición}        & El usuario ve datos actualizados de su cartera. \\
		\textbf{Excepciones}          & 
		\begin{itemize}
			\tightlist
			\item No hay valores posteriores en cartera (se muestran tablas vacías).
		\end{itemize} \\
		\textbf{Importancia}          & Alta \\
		\bottomrule
	\end{tabularx}
	\caption{CU-14 Eliminar valor de cartera.}
\end{table}


\begin{table}[p]
	\centering
	\begin{tabularx}{\linewidth}{ p{0.21\columnwidth} p{0.71\columnwidth} }
		\toprule
		\textbf{CU-15}    & \textbf{Crear un nuevo valor en seguimiento}\\
		\toprule
		\textbf{Versión}              & 1.0    \\
		\textbf{Autor}                & Rodrigo Merino Tovar \\
		\textbf{Requisitos asociados} & RF-3 y RF-3.3 \\
		\textbf{Descripción}          & Guardar un nuevo valor para realizar seguimiento.\\
		\textbf{Precondición}         & El usuario está logueado, las bases de datos están disponibles y el usuario no tiene el valor previamente en seguimiento.  \\
		\textbf{Acciones}             &
		\begin{enumerate}
			\def\labelenumi{\arabic{enumi}.}
			\tightlist
			\item El usuario accede al \emph{DashBoard}. 
			\item Mostrar información relativa a cartera.
			\item El usuario pulsa el botón \emph{Nuevo valor a seguir}.  
			\item Mostrar formulario de búsqueda de valores (\emph{tickers}).
			\item Seleccionar valor a seguir.
			\item Pulsar \emph{Guardar}.
			\item Buscar valores del mismo sector.
			\item Volver al \emph{DashBoard}
		\end{enumerate}\\
		\textbf{Postcondición}        & El usuario ve nuevo \emph{stock} en seguimiento junto con aquellos valores que pertenezcan al mismo sector. \\
		\textbf{Excepciones}          & 
		\begin{itemize}
			\tightlist
			\item No hay valores posteriores en seguimiento (se muestran tablas vacías).
		\end{itemize} \\
		\textbf{Importancia}          & Media \\
		\bottomrule
	\end{tabularx}
	\caption{CU-15 Crear un nuevo valor en seguimiento.}
\end{table}


\begin{table}[p]
	\centering
	\begin{tabularx}{\linewidth}{ p{0.21\columnwidth} p{0.71\columnwidth} }
		\toprule
		\textbf{CU-16}    & \textbf{Eliminar un valor de seguimiento}\\
		\toprule
		\textbf{Versión}              & 1.0    \\
		\textbf{Autor}                & Rodrigo Merino Tovar \\
		\textbf{Requisitos asociados} & RF-3 y RF-3.4 \\
		\textbf{Descripción}          & Para eliminar un valor que estuviera en seguimiento.\\
		\textbf{Precondición}         & El usuario está logueado, las bases de datos están disponibles y el usuario tiene previamente valores en seguimiento.  \\
		\textbf{Acciones}             &
		\begin{enumerate}
			\def\labelenumi{\arabic{enumi}.}
			\tightlist
			\item El usuario accede al \emph{DashBoard}. 
			\item Mostrar información relativa a cartera.
			\item El usuario pulsa el botón \emph{Eliminar valor en seguimiento}.  
			\item Mostrar valores en seguimiento).
			\item Seleccionar valores a eliminar.
			\item Pulsar \emph{Eliminar seleccionados}.
			\item Volver al \emph{DashBoard}
		\end{enumerate}\\
		\textbf{Postcondición}        & El usuario ve valores en seguimiento actualizados. \\
		\textbf{Excepciones}          & 
		\begin{itemize}
			\tightlist
			\item No hay valores previos en seguimiento (se muestran tablas vacías).
			\item No hay valores posteriores en seguimiento (se muestran tablas vacías).
		\end{itemize} \\
		\textbf{Importancia}          & Media \\
		\bottomrule
	\end{tabularx}
	\caption{CU-16 Eliminar un valor de seguimiento.}
\end{table}


\begin{table}[p]
	\centering
	\begin{tabularx}{\linewidth}{ p{0.21\columnwidth} p{0.71\columnwidth} }
		\toprule
		\textbf{CU-17}    & \textbf{Consultar índice}\\
		\toprule
		\textbf{Versión}              & 1.0    \\
		\textbf{Autor}                & Rodrigo Merino Tovar \\
		\textbf{Requisitos asociados} & RF-4, RF-4.1, RF-4.2, RF-4.3, RF-4.4 y RF-1.2.1  \\
		\textbf{Descripción}          & Para mostrar datos de valores de un índice de forma agregada.\\
		\textbf{Precondición}         & Las bases de datos están disponibles.  \\
		\textbf{Acciones}             &
		\begin{enumerate}
			\def\labelenumi{\arabic{enumi}.}
			\tightlist
			\item El usuario selecciona el índice que quiere consultar. 
			\item Mostrar tabla con componentes del índice.
			\item Visualizar gráfica de evolución del índice en su conjunto.  
			\item Facilitar enlaces RSS a noticias relacionadas.
		\end{enumerate}\\
		\textbf{Postcondición}        & El usuario ve información del índice. \\
		\textbf{Excepciones}          & 
		\begin{itemize}
			\tightlist
			\item No hay datos de índice (mensaje de página inexistente).
		\end{itemize} \\
		\textbf{Importancia}          & Alta \\
		\bottomrule
	\end{tabularx}
	\caption{CU-17 Consultar índice.}
\end{table}


\begin{table}[p]
	\centering
	\begin{tabularx}{\linewidth}{ p{0.21\columnwidth} p{0.71\columnwidth} }
		\toprule
		\textbf{CU-18}    & \textbf{Mostrar tabla de valores}\\
		\toprule
		\textbf{Versión}              & 1.0    \\
		\textbf{Autor}                & Rodrigo Merino Tovar \\
		\textbf{Requisitos asociados} & RF-4 y RF-4.1  \\
		\textbf{Descripción}          & Para mostrar tabla de componentes de un índice.\\
		\textbf{Precondición}         & Las bases de datos están disponibles.  \\
		\textbf{Acciones}             &
		\begin{enumerate}
			\def\labelenumi{\arabic{enumi}.}
			\tightlist
			\item El usuario selecciona el índice que quiere consultar. 
			\item Mostrar tabla con componentes del índice.
			\item El usuario ordena según variación diaria.
		\end{enumerate}\\
		\textbf{Postcondición}        & El usuario consulta precios de todos los valores del índice. \\
		\textbf{Excepciones}          & 
		\begin{itemize}
			\tightlist
			\item No hay datos de índice (mensaje de página inexistente).
		\end{itemize} \\
		\textbf{Importancia}          & Alta \\
		\bottomrule
	\end{tabularx}
	\caption{CU-18 Mostrar tabla de valores.}
\end{table}


\begin{table}[p]
	\centering
	\begin{tabularx}{\linewidth}{ p{0.21\columnwidth} p{0.71\columnwidth} }
		\toprule
		\textbf{CU-19}    & \textbf{Consultar un valor del índice}\\
		\toprule
		\textbf{Versión}              & 1.0    \\
		\textbf{Autor}                & Rodrigo Merino Tovar \\
		\textbf{Requisitos asociados} & RF-4, RF-4.1, RF-4.2, RF-4.2.1, RF-4.2.2 y RF-4.2.3, RF-4.2.4, RF-4.2.5 y RF-4.2.6 \\
		\textbf{Descripción}          & Para consultar información detallada de un único valor cotizado.\\
		\textbf{Precondición}         & El usuario está logueado, las bases de datos están disponibles y el valor existe.  \\
		\textbf{Acciones}             &
		\begin{enumerate}
			\def\labelenumi{\arabic{enumi}.}
			\tightlist
			\item El usuario selecciona el índice del valor de su interés. 
			\item El usuario selecciona el valor que quiere consultar. 
			\item Dar acceso a gráfica interactiva.
			\item Mostrar distribución de retornos.
			\item Mostrar datos de último mes.
			\item Ver evolución del sector.
			\item Mostrar formulario de comparación con otros valores.
			\item Mostrar grafos de correlación.
		\end{enumerate}\\
		\textbf{Postcondición}        & El usuario ve toda la información ordenada en forma de tablas o gráficas; es posible interactuar (de forma básica) con la gráfica interactiva. \\
		\textbf{Excepciones}          & 
		\begin{itemize}
			\tightlist
			\item No hay datos del valor (mensaje de página inexistente).
		\end{itemize} \\
		\textbf{Importancia}          & Alta \\
		\bottomrule
	\end{tabularx}
	\caption{CU-19 Consultar un valor del índice.}
\end{table}


\begin{table}[p]
	\centering
	\begin{tabularx}{\linewidth}{ p{0.21\columnwidth} p{0.71\columnwidth} }
		\toprule
		\textbf{CU-20}    & \textbf{Mostrar gráfica interactiva}\\
		\toprule
		\textbf{Versión}              & 1.0    \\
		\textbf{Autor}                & Rodrigo Merino Tovar \\
		\textbf{Requisitos asociados} & RF-4, RF-4.1, RF-4.2 y RF-4.2.1 \\
		\textbf{Descripción}          & Para interactuar con temporalidades e indicadores de un valor cotizado.\\
		\textbf{Precondición}         & El usuario está logueado, las bases de datos están disponibles y el valor existe.  \\
		\textbf{Acciones}             &
		\begin{enumerate}
			\def\labelenumi{\arabic{enumi}.}
			\tightlist 
			\item El usuario selecciona el valor que quiere consultar. 
			\item Dar acceso a gráfica interactiva.
			\item Se varían temporalidades ajustando volumen y medias móviles.
			\item Se comprueban datos diarios.
		\end{enumerate}\\
		\textbf{Postcondición}        & El usuario interactúa con la gráfica del valor. \\
		\textbf{Excepciones}          & 
		\begin{itemize}
			\tightlist
			\item No hay datos del valor (mensaje de página inexistente).
		\end{itemize} \\
		\textbf{Importancia}          & Alta \\
		\bottomrule
	\end{tabularx}
	\caption{CU-20 Mostrar gráfica interactiva.}
\end{table}


\begin{table}[p]
	\centering
	\begin{tabularx}{\linewidth}{ p{0.21\columnwidth} p{0.71\columnwidth} }
		\toprule
		\textbf{CU-21}    & \textbf{Mostrar distribución de retornos}\\
		\toprule
		\textbf{Versión}              & 1.0    \\
		\textbf{Autor}                & Rodrigo Merino Tovar \\
		\textbf{Requisitos asociados} & RF-4, RF-4.1, RF-4.2 y RF-4.2.2 \\
		\textbf{Descripción}          & Para ver la distribución que siguen los retornos de un valor en el último año (252 sesiones).\\
		\textbf{Precondición}         & El usuario está logueado, las bases de datos están disponibles, el valor existe y tiene datos históricos suficientes.  \\
		\textbf{Acciones}             &
		\begin{enumerate}
			\def\labelenumi{\arabic{enumi}.}
			\tightlist 
			\item El usuario selecciona el valor que quiere consultar. 
			\item Dar acceso a gráfica de distribución de retornos.
		\end{enumerate}\\
		\textbf{Postcondición}        & El usuario consulta la distribución de los retornos de un valor. \\
		\textbf{Excepciones}          & 
		\begin{itemize}
			\tightlist
			\item No hay datos del valor (mensaje de página inexistente).
		\end{itemize} \\
		\textbf{Importancia}          & Baja \\
		\bottomrule
	\end{tabularx}
	\caption{CU-21 Mostrar distribución de retornos.}
\end{table}


\begin{table}[p]
	\centering
	\begin{tabularx}{\linewidth}{ p{0.21\columnwidth} p{0.71\columnwidth} }
		\toprule
		\textbf{CU-22}    & \textbf{Mostrar datos del último mes}\\
		\toprule
		\textbf{Versión}              & 1.0    \\
		\textbf{Autor}                & Rodrigo Merino Tovar \\
		\textbf{Requisitos asociados} & RF-4, RF-4.1, RF-4.2 y RF-4.2.3 \\
		\textbf{Descripción}          & Para ver tabla con precios de un valor en el último mes.\\
		\textbf{Precondición}         & El usuario está logueado, las bases de datos están disponibles, el valor existe y tiene datos históricos suficientes.  \\
		\textbf{Acciones}             &
		\begin{enumerate}
			\def\labelenumi{\arabic{enumi}.}
			\tightlist 
			\item El usuario selecciona el valor que quiere consultar. 
			\item Dar acceso a tabla con datos de último mes.
			\item Mostrar tabla con colores verde (alcista) y rojo (bajista).
		\end{enumerate}\\
		\textbf{Postcondición}        & El usuario consulta la tabla de precios del último mes. \\
		\textbf{Excepciones}          & 
		\begin{itemize}
			\tightlist
			\item No hay datos del valor (mensaje de página inexistente).
		\end{itemize} \\
		\textbf{Importancia}          & Media \\
		\bottomrule
	\end{tabularx}
	\caption{CU-22 Mostrar datos del último mes.}
\end{table}


\begin{table}[p]
	\centering
	\begin{tabularx}{\linewidth}{ p{0.21\columnwidth} p{0.71\columnwidth} }
		\toprule
		\textbf{CU-23}    & \textbf{Mostrar evolución del sector}\\
		\toprule
		\textbf{Versión}              & 1.0    \\
		\textbf{Autor}                & Rodrigo Merino Tovar \\
		\textbf{Requisitos asociados} & RF-4, RF-4.1, RF-4.2 y RF-4.2.4 \\
		\textbf{Descripción}          & Para ver comparación - en términos relativos - entre el valor y su sector de referencia.\\
		\textbf{Precondición}         & El usuario está logueado, las bases de datos están disponibles, el valor existe y tiene datos históricos suficientes.  \\
		\textbf{Acciones}             &
		\begin{enumerate}
			\def\labelenumi{\arabic{enumi}.}
			\tightlist 
			\item El usuario selecciona el valor que quiere consultar. 
			\item Obtener el sector de referencia del valor.
			\item Obtener listado de valores del mismo sector.
			\item Calcular media de evolución de todos los valores del mismo sector. 
			\item Ajustar datos de forma relativa (porcentual).
			\item Mostrar gráfica con evolución de valor y evolución de sector.
		\end{enumerate}\\
		\textbf{Postcondición}        & El usuario consulta la gráfica comparativa de valor con sector de referencia. \\
		\textbf{Excepciones}          & 
		\begin{itemize}
			\tightlist
			\item No hay datos del valor (mensaje de página inexistente).
			\item No hay datos de valores del sector (gráfica sin comparación).
		\end{itemize} \\
		\textbf{Importancia}          & Alta \\
		\bottomrule
	\end{tabularx}
	\caption{CU-23 Mostrar evolución del sector.}
\end{table}


\begin{table}[p]
	\centering
	\begin{tabularx}{\linewidth}{ p{0.21\columnwidth} p{0.71\columnwidth} }
		\toprule
		\textbf{CU-24}    & \textbf{Comparar con otros valores}\\
		\toprule
		\textbf{Versión}              & 1.0    \\
		\textbf{Autor}                & Rodrigo Merino Tovar \\
		\textbf{Requisitos asociados} & RF-4, RF-4.1, RF-4.2 y RF-4.2.5 \\
		\textbf{Descripción}          & Para ver comparación - en términos relativos - con otros valores seleccionados por el usuario.\\
		\textbf{Precondición}         & El usuario está logueado, las bases de datos están disponibles, el valor existe y tiene datos históricos suficientes.  \\
		\textbf{Acciones}             &
		\begin{enumerate}
			\def\labelenumi{\arabic{enumi}.}
			\tightlist 
			\item El usuario selecciona el valor que quiere consultar. 
			\item El usuario selecciona en un formulario otro valor con el que comparar. 
			\item Mostrar comparación relativa de precios de cierre. 
			\item Mostrar comparación relativa de rentabilidades diarias. 
		\end{enumerate}\\
		\textbf{Postcondición}        & Recargar página con nueva gráfica comparativa. \\
		\textbf{Excepciones}          & 
		\begin{itemize}
			\tightlist
			\item No hay datos del valor (mensaje de página inexistente).
			\item No hay datos del valor con el que comparar (mensaje de \emph{ticker} no válido).
		\end{itemize} \\
		\textbf{Importancia}          & Media \\
		\bottomrule
	\end{tabularx}
	\caption{CU-24 Comparar con otros valores.}
\end{table}


\begin{table}[p]
	\centering
	\begin{tabularx}{\linewidth}{ p{0.21\columnwidth} p{0.71\columnwidth} }
		\toprule
		\textbf{CU-25}    & \textbf{Mostrar grafos de correlación}\\
		\toprule
		\textbf{Versión}              & 1.0    \\
		\textbf{Autor}                & Rodrigo Merino Tovar \\
		\textbf{Requisitos asociados} & RF-4, RF-4.1, RF-4.2 y RF-4.2.5 \\
		\textbf{Descripción}          & Para ver grafos de correlación positiva y negativa con el resto de valores disponibles.\\
		\textbf{Precondición}         & El usuario está logueado, las bases de datos están disponibles, el valor existe y tiene datos históricos suficientes.  \\
		\textbf{Acciones}             &
		\begin{enumerate}
			\def\labelenumi{\arabic{enumi}.}
			\tightlist 
			\item El usuario selecciona el valor que quiere consultar. 
			\item Crear matriz de correlación entre todos los valores disponibles.
			\item Crear grafo de correlación positiva (correl. mayor a 0,75)
			\item Crear grafo de correlación negativa (correl. mayor a -0,75)
			\item Mostrar grafos de correlaciones.
		\end{enumerate}\\
		\textbf{Postcondición}        & El usuario consulta los grafos de correlaciones. \\
		\textbf{Excepciones}          & 
		\begin{itemize}
			\tightlist
			\item No hay datos del valor (mensaje de página inexistente).
			\item No hay correlaciones que superen los umbrales (mostrar grafo con nodo de valor sin más datos).
		\end{itemize} \\
		\textbf{Importancia}          & Alta \\
		\bottomrule
	\end{tabularx}
	\caption{CU-25 Mostrar grafos de correlación.}
\end{table}


\begin{table}[p]
	\centering
	\begin{tabularx}{\linewidth}{ p{0.21\columnwidth} p{0.71\columnwidth} }
		\toprule
		\textbf{CU-26}    & \textbf{Mostrar noticias relacionadas}\\
		\toprule
		\textbf{Versión}              & 1.0    \\
		\textbf{Autor}                & Rodrigo Merino Tovar \\
		\textbf{Requisitos asociados} & RF-4 y RF-4.3 \\
		\textbf{Descripción}          & Para facilitar enlaces de fuentes RSS que estén relacionadas con el índice.\\
		\textbf{Precondición}         & Las bases de datos están disponibles y las fuentes RSS están activas.  \\
		\textbf{Acciones}             &
		\begin{enumerate}
			\def\labelenumi{\arabic{enumi}.}
			\tightlist 
			\item El usuario selecciona el índice que quiere consultar. 
			\item Mostrar listado de noticias relacionadas de fuentes RSS.
			\item Pulsar el botón \emph{leer más} de una fuente.
			\item Redirigir a web externa. 
		\end{enumerate}\\
		\textbf{Postcondición}        & El usuario ve listado de noticias relacionadas relevantes y puede seleccionar \emph{leer más} en fuente externa. \\
		\textbf{Excepciones}          & 
		\begin{itemize}
			\tightlist
			\item No hay datos de fuentes RSS (mensaje de error).
		\end{itemize} \\
		\textbf{Importancia}          & Baja \\
		\bottomrule
	\end{tabularx}
	\caption{CU-26 Mostrar noticias relacionadas.}
\end{table}


\begin{table}[p]
	\centering
	\begin{tabularx}{\linewidth}{ p{0.21\columnwidth} p{0.71\columnwidth} }
		\toprule
		\textbf{CU-27}    & \textbf{Gestionar \emph{Lab}}\\
		\toprule
		\textbf{Versión}              & 1.0    \\
		\textbf{Autor}                & Rodrigo Merino Tovar \\
		\textbf{Requisitos asociados} & RF-5, RF-5.1, RF-5.1.1, RF-5.1.2, RF-5.1.3, RF-5.1.4, RF-5.2, RF-5.2.1 y RF-5.2.2 \\
		\textbf{Descripción}          & Para controlar una sección de \emph{laboratorio} virtual con el que experimentar sobre los datos disponibles.\\
		\textbf{Precondición}         & El usuario está logueado y las bases de datos están disponibles.  \\
		\textbf{Acciones}             &
		\begin{enumerate}
			\def\labelenumi{\arabic{enumi}.}
			\tightlist 
			\item El usuario pulsa en en el apartado de \emph{Lab}.
			\item Se muestran opciones de modelado ARIMA y de \emph{trading} algorítmico.
			\item Para los usuarios del repositorio en local se muestra una opción de \emph{forecasting} con redes LSTM.
			\item Al pulsar cualquier opción disponible se facilitan los accesos a las funcionalidades concretas.
			\item Al lado de las funcionalidades disponibles se muestra información de ayuda al usuario.  
		\end{enumerate}\\
		\textbf{Postcondición}        & El usuario puede consultar las herramientas disponibles dentro del \emph{Lab}. \\
		\textbf{Excepciones}          & \\
		\textbf{Importancia}          & Alta \\
		\bottomrule
	\end{tabularx}
	\caption{CU-27 Gestionar \emph{Lab}.}
\end{table}


\begin{table}[p]
	\centering
	\begin{tabularx}{\linewidth}{ p{0.21\columnwidth} p{0.71\columnwidth} }
		\toprule
		\textbf{CU-28}    & \textbf{Trabajar con ARIMA}\\
		\toprule
		\textbf{Versión}              & 1.0    \\
		\textbf{Autor}                & Rodrigo Merino Tovar \\
		\textbf{Requisitos asociados} & RF-5, RF-5.1, RF-5.1.1, RF-5.1.2, RF-5.1.3 y RF-5.1.4  \\
		\textbf{Descripción}          & Para tratar de hacer \emph{forecasting} de series temporales con modelos ARIMA .\\
		\textbf{Precondición}         & El usuario está logueado y las bases de datos están disponibles.  \\
		\textbf{Acciones}             &
		\begin{enumerate}
			\def\labelenumi{\arabic{enumi}.}
			\tightlist 
			\item El usuario entra en el \emph{Lab}.
			\item Se pulsa el botón \emph{ARIMA}.
			\item Se muestran las opciones para trabajar con modelos ARIMA. 
			\item Se despliegan los botones de las herramientas disponibles: funciones ACF y PACF, modelado manual, modelado automático y modelado con búsqueda de parámetros por rejilla. 
			\item Al lado de cada opción se muestra una breve ayuda. 
			\item Pulsar de la opción deseada.
			\item Redirigir a la web con el formulario correpondiente. 
		\end{enumerate}\\
		\textbf{Postcondición}        & El usuario es dirigido a la plantilla con el formulario adecuado. \\
		\textbf{Excepciones}          & 
		\begin{itemize}
			\tightlist
			\item La redirección no es adecuada (mensaje de error inesperado).
		\end{itemize} \\
		\textbf{Importancia}          & Alta \\
		\bottomrule
	\end{tabularx}
	\caption{CU-28 Trabajar con ARIMA.}
\end{table}


\begin{table}[p]
	\centering
	\begin{tabularx}{\linewidth}{ p{0.21\columnwidth} p{0.71\columnwidth} }
		\toprule
		\textbf{CU-29}    & \textbf{Buscar (p,d,q) con funciones ACF y PACF:}\\
		\toprule
		\textbf{Versión}              & 1.0    \\
		\textbf{Autor}                & Rodrigo Merino Tovar \\
		\textbf{Requisitos asociados} & RF-5, RF-5.1 y RF-5.1.1  \\
		\textbf{Descripción}          & Para buscar los parámetros (p,d,q) de un modelo ARIMA de forma visual, a través del uso de funciones ACF y PACF.\\
		\textbf{Precondición}         & El usuario está logueado, las bases de datos están disponibles, el valor buscado existe.  \\
		\textbf{Acciones}             &
		\begin{enumerate}
			\def\labelenumi{\arabic{enumi}.}
			\tightlist 
			\item El usuario accede a la sección de ARIMA dentro del \emph{Lab}.
			\item Se pulsa el botón \emph{Buscar parámetros (p,d,q) con funciones ACF y PACF}. 
			\item Se accede al formulario de interacción adecuado. 
			\item El usuario introduce el valor buscado y la cantidad de sesiones previas que desea utilizar para el modelo. 
			\item El usuario pulsa el botón \emph{Mostrar resultados}.
			\item Se muestran las diferenciaciones más habituales: d=1, d=2 y d=3.			
			\item Se muestran las gráficas de ACF y PACF junto con precios de cierre diferenciados.
		\end{enumerate}\\
		\textbf{Postcondición}        & El usuario puede consultar la información y sacar conclusiones de los mejores parámetros para la serie temporal escogida. \\
		\textbf{Excepciones}          & 
		\begin{itemize}
			\tightlist
			\item Se indica un valor que no existe (mensaje de error y reintento).
			\item Se indica una cantidad inadecuada de sesiones (mensaje de error y reintento).
		\end{itemize} \\
		\textbf{Importancia}          & Alta \\
		\bottomrule
	\end{tabularx}
	\caption{CU-29 Buscar (p,d,q) con funciones ACF y PACF.}
\end{table}


\begin{table}[p]
	\centering
	\begin{tabularx}{\linewidth}{ p{0.21\columnwidth} p{0.71\columnwidth} }
		\toprule
		\textbf{CU-30}    & \textbf{Introducir (p,d,q) manualmente:}\\
		\toprule
		\textbf{Versión}              & 1.0    \\
		\textbf{Autor}                & Rodrigo Merino Tovar \\
		\textbf{Requisitos asociados} & RF-5, RF-5.1 y RF-5.1.2  \\
		\textbf{Descripción}          & Para realizar \emph{forecasting} con un modelo ARIMA que tenga los parámetros (p,d,q) deseados por el usuario.\\
		\textbf{Precondición}         & El usuario está logueado, las bases de datos están disponibles, el valor buscado existe y el usuario conoce los parámetros que va a utilizar.  \\
		\textbf{Acciones}             &
		\begin{enumerate}
			\def\labelenumi{\arabic{enumi}.}
			\tightlist 
			\item El usuario accede a la sección del \emph{Lab}.
			\item Se pulsa el botón \emph{Parámetros (p,d,q) introducidos por usuario}. 
			\item Se accede al formulario de interacción adecuado. 
			\item El usuario introduce valor buscado y cantidad de sesiones que desea utilizar con el modelo. 
			\item El usuario introduce la cantidad de datos que quiere utilizar para entrenar el modelo.
			\item El usuario introduce los parámetros (p,d,q).
			\item El usuario pulsa el botón \emph{Estimar}.
			\item Calcular intervalos de confianza y validación \emph{WF}.
			\item Calcular aciertos de tendencia al usar validación \emph{WF}.
			\item Se hace predicción de ARIMA sobre el porcentaje de datos de entrenamiento.
			\item Se muestran métricas de error: MSE y RMSE. 
			\item Se muestra comparación con una predicción naíf.			
			\item Se muestra gráfica con validación \emph{walk forward}.
			\item Mostrar resultados sobre datos de test con validación \emph{WF}.
		\end{enumerate}\\
		\textbf{Postcondición}        & El usuario puede consultar la información. \\
		\textbf{Excepciones}          & 
		\begin{itemize}
			\tightlist
			\item Se indica un valor que no existe (mensaje de error y reintento).
			\item Se indica una cantidad inadecuada de sesiones (mensaje de error y reintento).
			\item Se indica un porcentaje de datos de entrenamiento no válido (mensaje de error y reintento).
			\item Se indican parámetros (p,d,q) no válidos (mensaje de error y reintento).
		\end{itemize} \\
		\textbf{Importancia}          & Alta \\
		\bottomrule
	\end{tabularx}
	\caption{CU-30 Introducir (p,d,q) manualmente.}
\end{table}


\begin{table}[p]
	\centering
	\begin{tabularx}{\linewidth}{ p{0.21\columnwidth} p{0.71\columnwidth} }
		\toprule
		\textbf{CU-31}    & \textbf{Estimar (p,d,q) de forma automática:}\\
		\toprule
		\textbf{Versión}              & 1.0    \\
		\textbf{Autor}                & Rodrigo Merino Tovar \\
		\textbf{Requisitos asociados} & RF-5, RF-5.1 y RF-5.1.3  \\
		\textbf{Descripción}          & Para realizar \emph{forecasting} con un modelo ARIMA que tenga los parámetros (p,d,q) calculados de forma automática.\\
		\textbf{Precondición}         & El usuario está logueado, las bases de datos están disponibles y el valor buscado existe.  \\
		\textbf{Acciones}             &
		\begin{enumerate}
			\def\labelenumi{\arabic{enumi}.}
			\tightlist 
			\item El usuario accede a la sección de ARIMA dentro del \emph{Lab}.
			\item Se pulsa el botón \emph{Parámetros (p,d,q) estimados automáticamente}. 
			\item Se accede al formulario de interacción adecuado. 
			\item El usuario introduce el valor buscado y la cantidad de sesiones previas que desea utilizar para el modelo. 
			\item El usuario introduce la cantidad de datos que quiere utilizar para entrenar el modelo.
			\item El usuario pulsa el botón \emph{Estimar}.
			\item Se utiliza \emph{auto\_arima} para estimar parámetros.
			\item Calcular intervalos de confianza y validación \emph{WF}.
			\item Calcular aciertos de tendencia al usar \emph{walk forward}
			\item Se hace predicción de ARIMA sobre el porcentaje de datos de entrenamiento.
			\item Se muestran métricas de error: MSE y RMSE. 
			\item Se muestra comparación con una predicción naíf.			
			\item Se muestra gráfica con validación \emph{walk forward}.
			\item Mostrar resultados sobre datos de test con \emph{WF}.
			\item Se muestra informe ARIMA/SARIMAX.
		\end{enumerate}\\
		\textbf{Postcondición}        & El usuario puede consultar la información. \\
		\textbf{Excepciones}          & 
		\begin{itemize}
			\tightlist
			\item Se indica un valor que no existe (mensaje de error y reintento).
			\item Se indica una cantidad inadecuada de sesiones (mensaje de error y reintento).
			\item Se indica un porcentaje de datos de entrenamiento no válido (mensaje de error y reintento).
		\end{itemize} \\
		\textbf{Importancia}          & Alta \\
		\bottomrule
	\end{tabularx}
	\caption{CU-31 Estimar (p,d,q) de forma automática.}
\end{table}


\begin{table}[p]
	\centering
	\begin{tabularx}{\linewidth}{ p{0.21\columnwidth} p{0.71\columnwidth} }
		\toprule
		\textbf{CU-32}    & \textbf{Hacer búsqueda por rejilla para (p,d,q):}\\
		\toprule
		\textbf{Versión}              & 1.0    \\
		\textbf{Autor}                & Rodrigo Merino Tovar \\
		\textbf{Requisitos asociados} & RF-5, RF-5.1 y RF-5.1.4  \\
		\textbf{Descripción}          & Para realizar \emph{forecasting} con un modelo ARIMA que tenga los parámetros (p,d,q) obtenidos tras una búsqueda por rejilla de entre unos valores preestablecidos.\\
		\textbf{Precondición}         & El usuario está logueado, las bases de datos están disponibles, el valor buscado existe, los posibles valores de los parámetros están preestablecidos.  \\
		\textbf{Acciones}             &
		\begin{enumerate}
			\def\labelenumi{\arabic{enumi}.}
			\tightlist 
			\item El usuario accede a la sección de ARIMA.
			\item Se pulsa el botón \emph{Parámetros (p,d,q) creados con búsqueda por rejilla}. 
			\item Se accede al formulario de interacción adecuado. 
			\item El usuario introduce el valor buscado y la cantidad de sesiones previas que desea utilizar. 
			\item El usuario selecciona posibles valores (p,d,q).
			\item El usuario pulsa el botón \emph{Estimar}.
			\item Se evalúan todas las posibles combinaciones de parámetros y se selecciona el modelo con menor MSE (es decir, el mejor modelo posible). 
			\item Calcular intervalos de confianza y validación \emph{walk forward}.
			\item Se calculan aciertos de tendencia al usar \emph{WF}.
			\item Se hace predicción de ARIMA sobre el porcentaje de datos de entrenamiento.
			\item Se muestran métricas de error: MSE y RMSE. 
			\item Se muestra comparación con una predicción naíf.			
			\item Se muestra gráfica con validación \emph{walk forward}.
			\item Se muestran resultados sobre datos de test con \emph{walk forward} e informe ARIMA/SARIMAX.
		\end{enumerate}\\
		\textbf{Postcondición}        & El usuario puede consultar la información. \\
		\textbf{Excepciones}          & 
		\begin{itemize}
			\tightlist
			\item Se indica un valor que no existe (mensaje de error y reintento).
			\item Se indica una cantidad inadecuada de sesiones (mensaje de error y reintento).
			\item Se indica un porcentaje de datos de entrenamiento no válido (mensaje de error y reintento).
			\item Se indican posibles valores de (p,d,q) no válidos (mensaje de error y reintento).
		\end{itemize} \\
		\textbf{Importancia}          & Media \\
		\bottomrule
	\end{tabularx}
	\caption{CU-32 Hacer búsqueda por rejilla para (p,d,q).}
\end{table}


\begin{table}[p]
	\centering
	\begin{tabularx}{\linewidth}{ p{0.21\columnwidth} p{0.71\columnwidth} }
		\toprule
		\textbf{CU-33}    & \textbf{Trabajar con \emph{trading} algorítmico:}\\
		\toprule
		\textbf{Versión}              & 1.0    \\
		\textbf{Autor}                & Rodrigo Merino Tovar \\
		\textbf{Requisitos asociados} & RF-5, RF-5.2, RF-5.2.1 y RF-5.2.2   \\
		\textbf{Descripción}          & Para poder utilizar técnicas de \emph{trading} algorítmico aplicadas a los valores disponibles.\\
		\textbf{Precondición}         & El usuario está logueado, las bases de datos están disponibles y el valor buscado existe.  \\
		\textbf{Acciones}             &
		\begin{enumerate}
			\def\labelenumi{\arabic{enumi}.}
			\tightlist 
			\item El usuario entra en el \emph{Lab}.
			\item Se pulsa el botón \emph{Trading algorítmico}.
			\item Se muestran las opciones para trabajar con técnicas de \emph{trading} algorítmico. 
			\item Se despliegan los botones de las herramientas disponibles: algoritmo de cruce de medias y estrategia basada en \emph{machine learning}.
			\item Al lado de cada opción se muestra una breve ayuda. 
			\item Pulsar de la opción deseada.
			\item Redirigir a la web con el formulario correpondiente. 
		\end{enumerate}\\
		\textbf{Postcondición}        & El usuario puede consultar la información y sacar conclusiones sobre los resultados tanto en gráficas como en resumen descriptivo del modelo. \\
		\textbf{Excepciones}          & 
		\begin{itemize}
			\tightlist
			\item La redirección no es adecuada (mensaje de error inesperado).
		\end{itemize} \\
		\textbf{Importancia}          & Alta \\
		\bottomrule
	\end{tabularx}
	\caption{CU-33 Trabajar con \emph{trading} algorítmico.}
\end{table}


\begin{table}[p]
	\centering
	\begin{tabularx}{\linewidth}{ p{0.21\columnwidth} p{0.71\columnwidth} }
		\toprule
		\textbf{CU-34}    & \textbf{Usar algoritmo de cruce de medias:}\\
		\toprule
		\textbf{Versión}              & 1.0    \\
		\textbf{Autor}                & Rodrigo Merino Tovar \\
		\textbf{Requisitos asociados} & RF-5, RF-5.2 y RF-5.2.1   \\
		\textbf{Descripción}          & Para utilizar el algoritmo de cruce de medias sobre un valor seleccionado.\\
		\textbf{Precondición}         & El usuario está logueado, las bases de datos están disponibles y el valor buscado existe.  \\
		\textbf{Acciones}             &
		\begin{enumerate}
			\def\labelenumi{\arabic{enumi}.}
			\tightlist 
			\item El usuario accede a la sección de \emph{trading} algorítmico dentro del \emph{Lab}.
			\item Se pulsa el botón \emph{Algoritmo de cruce de medias}. 
			\item Se accede al formulario de interacción adecuado. 
			\item El usuario introduce el valor buscado y la cantidad de sesiones previas que desea utilizar para el algoritmo.
			\item Se hace una búsqueda por rejilla calculando las mejores rentabilidades con distintos cruces de medias preestablecidos ([30, 50, 200] y [10, 15, 20, 50]).
			\item Se muestra el mejor resultado posible. 
			\item Se hace una comparación con una estrategia \emph{buy and hold} en el mismo número de sesiones seleccionado.
		\end{enumerate}\\
		\textbf{Postcondición}        & El usuario puede ver un resultado porcentual de haber seguido la estrategia de cruce de medias en el período seleccionado y puede comparar con una estrategia \emph{buy and hold} para sacar conclusiones. \\
		\textbf{Excepciones}          & 
		\begin{itemize}
			\tightlist
			\item Se indica un valor que no existe (mensaje de error y reintento).
			\item Se indica una cantidad inadecuada de sesiones (mensaje de error y reintento).
		\end{itemize} \\
		\textbf{Importancia}          & Alta \\
		\bottomrule
	\end{tabularx}
	\caption{CU-34 Usar algoritmo de cruce de medias.}
\end{table}


\begin{table}[p]
	\centering
	\begin{tabularx}{\linewidth}{ p{0.21\columnwidth} p{0.71\columnwidth} }
		\toprule
		\textbf{CU-35}    & \textbf{Usar estrategia basada en \emph{machine learning}:}\\
		\toprule
		\textbf{Versión}              & 1.0    \\
		\textbf{Autor}                & Rodrigo Merino Tovar \\
		\textbf{Requisitos asociados} & RF-5, RF-5.2 y RF-5.2.2   \\
		\textbf{Descripción}          & Para hacer una predicción para la siguiente sesión de un índice con modelos de regresión y de clasificación.\\
		\textbf{Precondición}         & El usuario está logueado, las bases de datos están disponibles y el índice sobre el que predecir existe.  \\
		\textbf{Acciones}             &
		\begin{enumerate}
			\def\labelenumi{\arabic{enumi}.}
			\tightlist 
			\item El usuario accede a la sección de \emph{trading} algorítmico dentro del \emph{Lab}.
			\item Se pulsa el botón \emph{Estrategia basada en ML}. 
			\item Se accede al formulario de interacción adecuado. 
			\item El usuario introduce el valor buscado y la cantidad de sesiones previas que desea utilizar para el algoritmo.
			\item El usuario indica el tipo de modelo. 
			\item El usuario indica el porcentaje de datos para entrenar el modelo.
			\item Se seleccionan los datos de los valores que representan al índice para el mismo período (mismo número de sesiones).
			\item Se entrena el modelo y se calculan los \emph{scores} en entrenamiento y en test. 
			\item Se hace una predicción adicional, fuera de los datos de test. 
			\item Se muestra al usuario información tabular con resultados obtenidos y con predicción para la siguiente sesión. 
		\end{enumerate}\\
		\textbf{Postcondición}        & El usuario puede consultar una tabla de resultados tras entrenar el modelo y la predicción. \\
		\textbf{Excepciones}          & 
		\begin{itemize}
			\tightlist
			\item Se indica una cantidad inadecuada de sesiones (mensaje de error y reintento).
			\item Se indica un índice que no existe (mensaje de error y reintento).
			\item Se indica un tipo de modelo no disponible (mensaje de error y reintento).
			\item Se indica un porcentaje de datos de entrenamiento no válido (mensaje de error y reintento).
		\end{itemize} \\
		\textbf{Importancia}          & Alta \\
		\bottomrule
	\end{tabularx}
	\caption{CU-35 Usar estrategia basada en \emph{machine learning}.}
\end{table}
\apendice{Especificación de diseño}

\section{Introducción}

\section{Diseño de datos}

\section{Diseño procedimental}

\section{Diseño arquitectónico}



\apendice{Documentación técnica de programación}

\section{Introducción}

En este anexo se recoge la documentación técnica de programación, que incluye recomendaciones para el entorno de desarrollo, la estructura de directorio, los procesos de compilación, la configuración de integración e instalación de dependencias y las baterías de tests realizadas. 

\section{Estructura de directorios}

El repositorio del proyecto tiene la siguiente distribución de directorios:

\begin{itemize}
\tightlist
\item

\end{itemize}
\section{Manual del programador}

\section{Compilación, instalación y ejecución del proyecto}

\section{Pruebas del sistema}

\apendice{Documentación de usuario}

\section{Introducción}

En esta sección se presenta un breve manual de usuario que permita interaccionar con la herramienta de manera satisfactoria. Aquí se muestran los principales formularios y las opciones que tiene el usuario final para trabajar con este entorno. 

\section{Requisitos de usuarios}

\subsection{Uso en web}

Este proyecto está especialmente pensado para ser utilizado en un entorno local, pero se ha dispuesto una página web para que un usuario pueda hacer pruebas antes de instalarlo en su equipo. Visitar \href{http://takeiteasy.pythonanywhere.com/}{FAT: Financial Analysis Tool}.

Los requisitos para usar esta funcionalidad no son más que disponer de un navegador web. De entre los más populares y modernos se ha probado la integración con \emph{Google Chrome}, \emph{Mozilla Firefox}, \emph{Brave Browser} y \emph{Microsoft Edge} con resultados satisfactorios. 

\subsection{Uso en entorno local}

Los requisitos previos para poder utilizar este proyecto se describen en los \textbf{pasos 1 a 3 que se indican en el manual del programador} \ref{manual_programador}.


\section{Instalación}

Para utilizar este proyecto en un entorno local lo mejor es seguir los \textbf{pasos 4 a 9 que se indican en el manual del programador} \ref{manual_programador}. 

\section{Manual del usuario}

\subsection{Página principal}

\imagen{img_anex_32_usuario_01.png}{Página principal}{1}

En portada están disponibles unos \emph{carruseles} de noticias que nos muestran información en inglés relacionadas con los principales mercados mundiales. Además, tenemos otros carruseles que permiten ver cuáles han sido los mejores y peores valores de cada índice en la última sesión. 

Desde aquí podemos consultar información detallada de un índice y podemos registrarnos o loguearnos (si ya estamos registrados).

\subsection{Registro y \emph{login}}

Para registrarnos se deberán cumplir unas condiciones mínimas de complejidad de contraseña y de nombre de usuario. En caso de no cumplir los condicionantes el usuario será informado:

\begin{itemize}
\tightlist
\item
Nombre: al menos 5 caracteres. Sólo se permiten números y letras.
\item
Contraseña: al menos 8 caracteres. No se permiten contraseñas \emph{habituales}. 
\end{itemize}

Una vez registrados se nos redirige a la página principal, ya logueados, y se nos abrirán las opciones de \emph{DashBoard} (área de usuario) y \emph{Lab} (laboratorio virtual para experimentar con series temporales):

\imagen{img_anex_33_usuario_02.png}{Proceso de registro}{1}

\newpage

Desde cualquier página de la herramienta tendremos disponible el \emph{login} y el registro si hemos accedido previamente. En cualquier caso, cada vez que nos logueamos, la web nos redirige al \emph{DashBoard} (área personal):

\imagen{img_anex_34_usuario_03.png}{Proceso de \emph{login}}{1}

\subsection{Consultar un índice bursátil}

Desde la página principal, desde el área de usuario, etc. podremos acceder a información sobre cualquiera de los índices disponibles. Sólo hay que hacer \emph{click} en los botones disponibles y se abrirá la ventana con los datos requeridos. 

Los datos que podremos consultar son:

\begin{itemize}
\tightlist
\item
Precios de apertura, cierre, máximo, mínimo y variación de cada uno de los componentes del índice.
\item
Gráfica con la evolución del índice en el último año (aproximadamente).
\item
Noticias relacionadas con el índice bursátil.
\end{itemize}

Las noticias se obtienen de \emph{feeds} RSS y es posibles añadir o modificar las fuentes disponibles. Para ello, si tienes la herramienta instalada en local, sólo tienes que modificar el archivo \texttt{/util/rss/RSS.py}. 

\imagen{img_anex_35_usuario_04.png}{Consultar índice bursátil}{1}


\subsection{Consultar un valor}

Para poder consultar un valor es necesario estar logueado. Lo habitual será consultar un valor haciendo \emph{click} desde las tablas de los índices. Una vez seleccionamos un valor se nos dirige a la ventana con sus datos, entre los que tenemos la gráfica interactiva:

\imagen{img_anex_36_usuario_05.png}{Consultar valor}{1}


También se puede consultar la evolución de precios del último mes y los grafos de correlación, que nos indican qué valores han evolucionado de manera similar en el último año y cuáles lo han hecho de manera \emph{inversa}:

\imagen{img_anex_37_usuario_06.png}{Consultar valor}{1}

Podemos ver más información, como la gráfica comparativa con la evolución del sector\footnote{Se calcula la evolución del sector considerando todos los valores, de todos los índices, disponibles en ese sector.} o la gráfica con la distribución de los retornos en porcentaje diario:

\imagen{img_anex_38_usuario_07.png}{Consultar valor}{1}

Por último, está disponible una sección de comparación con otros valores, para ver cómo han evolucionado, de forma relativa, en el último año:

\imagen{img_anex_39_usuario_08.png}{Consultar valor}{1}


\subsection{Utilizar el \emph{DashBoard}}

Para poder acceder al área personal es necesario estar logueados. 

\imagen{img_anex_40_usuario_09.png}{Acceso al \emph{DashBoard}}{0.9}

Como se aprecia en la imagen anterior, al acceder al \emph{DashBoard} la primera vez veremos unas tablas vacías, que nos indican que no tenemos valores en cartera ni valores en seguimiento.

Para añadir valores a nuestra cartera o eliminarlos sólo tenemos que utilizar los botones que aparecen en la parte superior del \emph{DashBoard}. El usuario debe tener en cuenta que los formularios de interacción con esta herramienta tienen valores limitados y, en muchos casos, evitan errores, lo que facilita el control de nuestras inversiones. Por ejemplo, si introducimos mal los datos de compra de un valor, la herramienta no nos lo permitirá:

\imagen{img_anex_41_usuario_10.png}{Acceso al \emph{DashBoard}}{1}

Según vamos añadiendo valores, nuestro \emph{DashBoard} va mostrando información adicional que puede ser de gran utilidad como, por ejemplo, la distribución de nuestra cartera - con los cálculos hechos tomando el euro como moneda de referencia - o la inversión que tenemos por sectores.

Y, por supuesto, podremos ver cómo está evolucionando nuestra cartera de forma individual o colectiva:

\imagen{img_anex_42_usuario_11.png}{Distribución de inversiones y sectores}{1}

\newpage
También podemos ver cómo está distribuida nuestra cartera por divisas:

\imagen{img_anex_43_usuario_12.png}{Distribución por divisas}{1}

Una de las partes más interesante es que podremos ver la gráfica que nos muestra la relación rentabilidad - riesgo que tenemos en ese momento. Es recomendable revisar el apartado teórico de la memoria adjunta a este proyecto, para comprender qué representan una gráfica de Markowitz y el ratio de Sharpe:


\imagen{img_anex_44_usuario_13.png}{Análisis de rentabilidad - riesgo}{1}

Es \textbf{importante} tener en cuenta que los valores utilizados para calcular la rentabilidad - riesgo son datos anuales, \textbf{no} se corresponden con los datos desde el momento de compra de los valores.


Para añadir o eliminar valores en seguimiento haremos un proceso similar al anterior:

\imagen{img_anex_45_usuario_14.png}{Añadir un valor en seguimiento}{1}

\newpage
\subsection{Utilizar el \emph{Lab}}

Para poder acceder al laboratorio virtual es necesario estar logueados y hacer \emph{click} en el botón de \emph{Lab} de la cabecera de la página. Esto nos llevará a la página principal del laboratorio virtual, desde donde podremos seleccionar la opción que más nos interesa para trabajar:

\imagen{img_anex_46_usuario_15.png}{Ventana principal del \emph{Lab}}{1}

En esta ventana podemos ver unos iconos de ayuda, que al pasar el ratón por encima, muestran información de utilidad para que el usuario sepa para qué sirve cada apartado. 

Si queremos trabajar, por ejemplo, con ARIMA y estamos pensando en buscar los mejores parámetros (p,d,q), podemos utilizar el apartado de las funciones ACF y PACF. La interpretación de estas gráficas queda fuera de las intenciones de este manual, pero se puede encontrar información relacionada en el apartado teórico de la memoria asociada de este trabajo:

\imagen{img_anex_47_usuario_16.png}{Usar ACF y PACF}{1}

Si lo que nos interesan es la funcionalidad que calculan los parámetros por nosotros (automáticamente o con búsqueda por rejilla) o queremos introducir unos datos para probar, seguiremos un criterio similar al anterior: 

\begin{itemize}
\item
Seleccionar la funcionalidad deseada.
\item
Rellenar el formulario con los datos necesarios. 
\item
Interpretar la información obtenida (en algunos apartados hay iconos de ayuda que puede aportar información adicional relevante).
\end{itemize}


En los formularios se solicita al usuario un porcentaje de entrenamiento, que corresponde a la cantidad de datos que se utilizarán para entrenar el modelo. En todos los casos las gráficas vienen acompañadas de unas tablas con los errores cometidos en los datos de test. Y se puede ver el resultado que obtendríamos si siguiéramos una estrategia de predicción naíf, que consiste en estimar que en la siguiente sesión el valor será el mismo que en la sesión actual (se puede consultar el icono de ayuda si existen dudas al respecto).

El usuario debe de tener en cuenta que los datos de test corresponden a los últimos datos de la serie temporal. Es decir, si tenemos una serie temporal de 500 sesiones y hemos dedicado el 90\% de los datos a entrenar el modelo, nos quedará un 10\% de datos de test y esos datos será las últimas 50 sesiones. Esto es así por el propio carácter de temporalidad en los datos analizados. 

\newpage
Por ejemplo, si queremos ver una predicción con un modelo ARIMA automático:

\imagen{img_anex_48_usuario_17.png}{Predicción automática de ARIMA}{1}

\newpage
Y al enviar el formulario se nos dirige a la página con los resultados:

\imagen{img_anex_49_usuario_18.png}{Resultados de ARIMA}{0.95}

Por otro lado, si lo que queremos es trabajar con funcionalidades de \emph{trading} algorítmico podemos escoger cualquiera de las dos opciones disponibles en el apartado correspondiente:

\begin{itemize}
\item
Algoritmo de cruce de medias.
\item
Estrategias basadas en \emph{machine learning}. 
\end{itemize}


Cuando se trabaja con el algoritmo de cruce de medias - con formularios similares a los de otras funcionalidades -, el usuario debe de tener en cuenta que la selección de las medias móviles simples se realiza de forma automática, buscando el mejor rendimiento pasado posible. Por tanto, no todos los valores analizados tendrán las mismas medias. Además, si se da la circunstancia de estar en estado de \emph{invertido} \textbf{en la última sesión disponible, se considerará que en esa sesión se cierra la posición para hacer los cálculos de rentabilidad}. Un ejemplo de resultado de este algoritmo es:

\imagen{img_anex_50_usuario_19.png}{Resultados algoritmo de cruce de medias}{0.9}

Por último, si queremos utilizar alguna de las estrategias basadas en \emph{machine learning}, con un modelo de regresión lineal o un modelo de clasificación, tendremos que seleccionar el botón correspondiente en el \emph{Lab}. Una vez seleccionada la opción deseada - de nuevo, a través de formularios que resultarán familiares al usuario - podremos observar los resultados obtenidos:

\imagen{img_anex_51_usuario_20.png}{Resultados estrategia basada en \emph{machine learning}}{1}


Para conocer más detalles sobre cómo se utiliza esta última estrategia se recomienda consultar la parte teórica de la memoria y, si el usuario desea, el método \texttt{estrategia\_machine\_learning()} en el archivo \texttt{/Lav/views.py} del código fuente. 
\apendice{Anexo de sostenibilización curricular}

\section{Introducción}

El presente Trabajo de Fin de Grado ha abarcado temas relacionados con \emph{Django}, finanzas, teoría moderna de porfolios o teoría de Harry M. Markowitz, Sharpe Ratio, \emph{machine learning} y ARIMA, entre otros. Durante la realización de este proyecto, se han adquirido y aplicado diversas competencias en sostenibilidad, según las directrices establecidas por la CRUE. En este anexo, se reflejarán las reflexiones personales sobre cómo los aspectos de sostenibilidad se han integrado y abordado a lo largo del trabajo.

\section{Competencias en sostenibilidad adquiridas}

\subsection{Economía sostenible}

Uno de los pilares fundamentales de la sostenibilidad es la economía. En el desarrollo de este trabajo, se ha aplicado la teoría de Harry M. Markowitz para la optimización de carteras y el Sharpe Ratio para evaluar el rendimiento ajustado al riesgo. Estas herramientas permiten a los inversores tomar decisiones más informadas y eficientes, minimizando riesgos y maximizando rendimientos, lo cual es esencial para una economía sostenible. Al aprender y aplicar estos conceptos, he adquirido una mayor comprensión de cómo las estrategias de inversión pueden contribuir a la estabilidad económica a largo plazo.

\subsection{Innovación y desarrollo tecnológico}

El uso de \emph{Django} como \emph{framework} para el desarrollo del proyecto destaca la importancia de la tecnología en la sostenibilidad. Este \emph{framework} es de código abierto, lo que promueve la colaboración y el intercambio de conocimiento, algo esencial para el desarrollo sostenible. Además, su eficiencia en el manejo de recursos y su escalabilidad aseguran que las aplicaciones desarrolladas sean sostenibles a largo plazo, tanto en términos de mantenimiento como de consumo de recursos.

\section{Impacto social y ambiental}

\subsection{Educación financiera}

La integración de conceptos financieros avanzados y técnicas de \emph{machine learning} en este trabajo no solo tienen un impacto positivo en términos de economía y tecnología, sino también en la educación financiera. Al compartir los resultados y metodologías del trabajo, se contribuye a mejorar el conocimiento y la conciencia financiera. Una sociedad con mayor educación financiera es capaz de tomar decisiones más sostenibles, tanto a nivel individual como colectivamente, promoviendo así un desarrollo económico equitativo y responsable.

\subsection{Licencia libre y repositorio público}

El proyecto está bajo una licencia \emph{CC BY-NC-SA 4.0} (Creative Commons Attribution-NonCommercial-ShareAlike 4.0 International) y está alojado en un repositorio público en \emph{GitHub}. Esto tiene varios impactos positivos en la sostenibilidad:

\begin{itemize}
\item
Promoción del conocimiento abierto.
\item
Colaboración global.
\item
Eficiencia en el uso de recursos tecnológico y humanos.
\end{itemize}

\subsection{Documentación del proyecto}

En este trabajo se ha dado especial relevancia a la documentación, tanto del código como de la teoría que fundamenta las decisiones tomadas durante el desarrollo. Esta documentación está disponible en \href{https://fat.readthedocs.io/es/latest/intro.html}{FAT. Read The Docs} y también se han creado \href{https://github.com/rmt0009alu/FAT/blob/main/docs/latex/anexos.pdf}{anexos} específicos para desarrolladores y usuarios. 

Este enfoque tiene los siguientes beneficios en términos de sostenibilidad:

\begin{itemize}
\item
Facilita la accesibilidad y la transparencia.
\item
Promueve el mantenimiento y la evolución del código. 
\item
Provee un recurso educativo para futuros estudiantes. 
\end{itemize}



\bibliographystyle{plain}
\bibliography{bibliografiaAnexos}

\end{document}
